\section{The one-loop term of the composite map}\label{sec:ol-composite}

This part prepares the ultraviolet limit of the four-dimensional theory: the growth $k\beta_\infty+O(1)$ of the one-loop coefficient of the running coupling summed over $k$ steps, in particular its positivity for large $k$. This is the input of our version, for the couplings of the small-field flow, of Ba\l aban's Theorem~2 \cite{BalYM4I}, which he states for $SU(2)$ and whose perturbative proof he deferred. This section reduces the one-loop coefficients of all steps to a single object, the one-loop term of the composite $k$-step map.

\paragraph{Why one coupling, and the plan of Part~XI.} The couplings $g_j$ of the flow change from step to step, but the object of this Part does not depend on them. The one-loop term $\Gamma^{\rm step}_j=-\frac12\log\det\Delta_j+\sum_c\log\mathfrak j_c$ of the step from scale $j$ is a function of the background field alone (paragraph ``The one-loop terms of the steps'' below), and the one-loop term $\Gamma^{(1)}_k$ of the composite map, obtained by Laplace's method as $g\to0$, does not depend on $g$ (Proposition~\ref{prop:ol-laplace}). So a single coupling $g$ serves as the parameter of Laplace's method. The couplings of the flow enter only in Section~\ref{sec:ol-uv}, where the coefficient $\beta_{j+1}$ of the flow is compared with the one-loop coefficient, through (F2) and (F3) of Proposition~\ref{prop:flow-F}. The chain is the following.
\begin{itemize}
\item Laplace's method defines $\Gamma^{(1)}_k$ (Proposition~\ref{prop:ol-laplace}); the integral over the intermediate scale gives the one-step relation (Lemma~\ref{lem:ol-one-step}), and induction gives $\Gamma^{(1)}_k=\sum_{j<k}\Gamma^{\rm step}_j$ up to a constant (Proposition~\ref{prop:ol-composition}).
\item The lattice symmetry gives the marginal tensor $\beta(\delta_{\mu\nu}|p|^2-p_\mu p_\nu)$ of each such function, with no term linear in $p$ (Lemma~\ref{lem:ol-symmetry}). In four dimensions the change of scale preserves marginal tensors (Lemma~\ref{lem:ol-pullback}), so the tensors of the steps add up to that of $\Gamma^{(1)}_k$ (Corollary~\ref{cor:ol-cumulative}). Section~\ref{subsec:ol-infinite-volume} builds the infinite-volume operators at $U=1$ that these statements and Sections~\ref{sec:ol-chart}--\ref{sec:ol-uv} use.
\item Section~\ref{sec:ol-chart} writes $\Gamma^{(1)}_k$ in the Landau chart; Sections~\ref{sec:ol-vacpol} and~\ref{sec:ol-feynman} compute its coefficient, $k\beta_\infty+O(1)$ (Theorem~\ref{thm:ol-cumulative}); Section~\ref{sec:ol-uv} transfers it to the flow (Proposition~\ref{prop:flow-F}) and proves Theorem~\ref{thm:ol-uv}.
\end{itemize}
The \hyperref[intro:XI]{introduction to Part~XI} describes the same chain.

\paragraph{What is used.} The section uses Parts~IV--VII: the one-step formula of Theorem~\ref{thm:step}, the exponential gauge fixing of Lemma~\ref{lem:exp-gauge}, and Lemma~\ref{lem:gf-step-invariance}(iv). It also uses Proposition~\ref{prop:jacobians}(c),(d), Proposition~\ref{prop:first-step} and Lemma~\ref{lem:first-step-formula} of Part~X, the latter two for the step from scale $0$; the lemma is the formula for the first step. Section~\ref{subsec:ol-infinite-volume} uses Part~I. A single coupling $g$ is kept throughout. This is legitimate in $d=4$, where the Wilson action of Section~\ref{sec:action} does not change under a change of units (Section~\ref{sec:rg-step} and the end of Theorem~\ref{thm:step}). The coupling of a step changes only through the extraction of $\beta$, which is not used here; the coupling paragraphs of Section~\ref{sec:rg-step} are not used either. Apart from the single coupling, the dimension enters the proofs of this section only through the scaling in Lemma~\ref{lem:ol-pullback}. The orientation in $d=4$ mentioned in Remark~\ref{rem:marginal-orders} is a comment: the proof of Lemma~\ref{lem:ol-symmetry} uses the point reflection $x\mapsto-x$, which lies in the hyperoctahedral group in every dimension.

\paragraph{Setting.} Fix $k\ge1$. A real configuration $V$ on $\T$ is \emph{small field} if:
\begin{itemize}
\item its curvature is so small that the hypotheses of Lemma~\ref{lem:composition} and of Theorem~\ref{thm:step} hold at every intermediate scale $j<k$, for the configurations $U_k(V)^{(j+1)}$ (at $j=0$, those of Lemma~\ref{lem:first-step-formula});
\item the side of the torus at every scale is at least the threshold $n_*$ of (T0), which contains $n_{\delta=1}$ and $80M'_L$.
\end{itemize}
In Part~XI ``small field'' always has this meaning: a condition on the curvature of $V$ and on the torus, which does not involve the coupling. It is not the small-field set $\mathfrak U_k(h_k)$ of Part~X, whose threshold $h_k=p(g_k)g_k$ depends on the coupling (Section~\ref{subsec:flow-parameters}).

\begin{readingbox}[Recall from Parts~IV and~V]
For a real reference $U_0$ put $W_k=U_0^{(k)}$, and let $\mathfrak B(V)=\{U:U^{(k)}=V\}$ be the constraint set. $\mathfrak U(\theta)$ is the set of real configurations with curvature and current at most $\theta$ (Section~\ref{sec:action}), and $\Phi(A)=S(e^{\ii\eta A}U_0)$. Corollary~\ref{cor:parametrization} parametrizes the orbits of $\mathfrak G_0$ in $\mathfrak B(V)\cap\mathfrak U(\theta)$ by fields in the Landau gauge: the set $\mathcal S_V$ of the real $A$ with $|A|_\infty\le r_u$, $RD^TA=0$, $e^{\ii\eta A}U_0\in\mathfrak U(\theta)$ and $\mathcal Q(e^{\ii\eta A})=VW_k^{-1}$ is mapped bijectively by $A\mapsto\mathsf U(A)$ onto the averaged axial representatives of these orbits (Definition~\ref{def:axial}), and $S(\mathsf U(A))=\Phi(A)$. Here $R$ is the projection $I-P_\cV$ of Part~II, $\mathcal Q$ the gauge-adapted averaging and $\mathsf w_k$, which enters $\mathsf U$, the correspondence transformation (Definition~\ref{def:adapted}). Below $U_0=U_k(V)$ lies in the minimal orbit $\mathscr M(V)$ of Theorem~\ref{thm:balaban}; then $W_k=V$, and the condition reads $\mathcal Q(e^{\ii\eta A})=1$.
\end{readingbox}
Let $\mathscr M(V)$ be the minimal orbit of $V$ (Theorem~\ref{thm:balaban}), $U_k(V)$ an element of it, and $\mathfrak G_0=\{u:u|_\T=1\}$. The group $\mathfrak G_0$ acts freely on all configurations: if $U^u=U$ and $u(y)=1$ at a point, parallel transport along any path from $y$ gives $u=1$.

Take $U_0=U_k(V)$ as the reference of Part~IV, and let $\theta\le\theta_0$ be such that $U_k(V)\in\mathfrak U(\theta')$ for some $\theta'<\theta$. The choice $\theta=\theta_0$ is allowed: by the hypotheses of Lemma~\ref{lem:composition} at $j=k-1$ (of Lemma~\ref{lem:first-step-formula} if $k=1$), $V\in\mathfrak U_k(\frac12a_1)$, so $U_k(V)\in\mathfrak U(\frac12B_3a_1)\subset\mathfrak U(\frac12\theta_0)$ by Theorem~\ref{thm:balaban}(b). The reference is $(M,B_3a_1)$-regular with $|J_{U_0}|_\infty\le B_3a_1$ and $U_0^{(k)}=V$, and $B_3a_1\le\frac12c_6$ (proof of Theorem~\ref{thm:balaban}, with $C^{\rm units}_L\ge2$ from Lemma~\ref{lem:units}); so Corollary~\ref{cor:parametrization} applies, with $\eps=B_3a_1$ and $\beta=0$ (the $\beta$ of that corollary, unrelated to the coefficients $\beta_j$ of the flow), since $\theta_0\le\frac13c_4$ and $c_6\le\frac13c_4$ (Section~\ref{subsec:reduction}). It gives a bijection $A\mapsto\mathsf U(A)$ from the set $\mathcal S_V$ onto the averaged axial representatives of the $\mathfrak G_0$-orbits in $\mathfrak B(V)\cap\mathfrak U(\theta)$, with $S(\mathsf U(A))=\Phi(A)$. Here $\mathcal S_V$ is the set of real $A$ with $|A|_\infty\le r_u$, $RD^TA=0$, $e^{\ii\eta A}U_0\in\mathfrak U(\theta)$ and $\mathcal Q(e^{\ii\eta A})=1$. For $r>0$ let $\mathscr N_r(V)$ be the union of the $\mathfrak G_0$-orbits of the $\mathsf U(A)$ with $A\in\mathcal S_V$, $|A|_\infty<r$. These definitions use only Theorem~\ref{thm:balaban} and Corollary~\ref{cor:parametrization}. So they make sense, by the same argument, for every real $V\in\mathfrak U_k(a_1)$, on a torus as in the setting, and every $\theta\le\theta_0$ with $\mathscr M(V)\subset\mathfrak U(\theta')$ for some $\theta'<\theta$; we call such a pair $(V,\theta)$ \emph{admissible}. By Lemma~\ref{lem:ol-chart-covariance}(b) below, $\mathscr N_r(V)$ does not depend on the choice of $U_k(V)$ in $\mathscr M(V)$. For $g>0$ put
\[
I_k(V;g,r)=\int\prod_bdU(b)\ \prod_c\delta\big(U^{(k)}(c)V(c)^{-1}\big)\ \one_{\mathscr N_r(V)}(U)\ e^{-S(U)/g^2},
\]
with the normalized Haar measure and the Haar $\delta$-function, as in \eqref{eq:rg-map}.

The neighbourhoods $\mathscr N_r(V)$ are defined through a chart whose reference is an element of the minimal orbit. The next lemma shows that they do not depend on this choice, that they are covariant, and that they shrink continuously to the minimal orbit. As in Section~\ref{subsec:var-symmetry}, a subscript or a second argument indicates the reference: $\mathcal S_V(U_0)$, $\mathsf U_{U_0}$, $\mathcal Q(\cdot\,;U_0)$, $\mathsf w_k(\cdot\,;U_0)$. For a real gauge transformation $u$ we write $R(u)$ for its action on Lie fields, $(R(u)A)(b)=R(u(b_-))A(b)$, as in Proposition~\ref{prop:linearized-averages}(f), and on scalar fields, $(R(u)\lambda)(x)=R(u(x))\lambda(x)$. To avoid a clash, in the proof below $R_{U_0}$ denotes the projection $R=I-P_\cV$ of Part~II at the reference $U_0$.

\begin{lemma}[The chart neighbourhoods]\label{lem:ol-chart-covariance}
\begin{enumerate}[label=(\alph*)]
\item \emph{(Gauge covariance of the chart.)} Let $U_0$, $V$, $\theta$ and $\beta$ be as in Corollary~\ref{cor:parametrization}, and let $u$ be a real gauge transformation on $\Tf$, with $v=u|_\T$. Then $U_0^u$ satisfies the hypotheses of Section~\ref{sec:chart} with the same $\eps$, $W_k(U_0^u)=W_k^v$, and $V^v$ satisfies the hypothesis of Corollary~\ref{cor:parametrization} relative to $U_0^u$, with the same $\beta$. The map $A\mapsto R(u)A$ maps $\mathcal S_V(U_0)$ onto $\mathcal S_{V^v}(U_0^u)$, preserves $|\cdot|_\infty$, and
\[
\mathsf U_{U_0^u}\big(R(u)A\big)=\mathsf U_{U_0}(A)^u\qquad(A\in\mathcal S_V(U_0)).
\]
\item \emph{(The choice of the reference.)} Let $(V,\theta)$ be admissible. Then $\mathscr N_r(V)$ does not depend on the choice of $U_k(V)$ in $\mathscr M(V)$. For every real gauge transformation $u$ on $\Tf$ the pair $(V^{u|_\T},\theta)$ is admissible, and $\mathscr N_r(V^{u|_\T})=\mathscr N_r(V)^u$.
\item \emph{(The chart radius.)} Let $\theta\le\theta_0$, and let $\mathcal K_0$ be a compact set of real configurations on $\T$ such that $(V,\theta)$ is admissible for every $V\in\mathcal K_0$. For $U\in\mathfrak U(\theta)$ with $U^{(k)}\in\mathcal K_0$ let $\varrho(U)=|A|_\infty$, where $\mathsf U(A)$, $A\in\mathcal S_{U^{(k)}}$, is the averaged axial representative of the $\mathfrak G_0$-orbit of $U$, for a reference in $\mathscr M(U^{(k)})$. Then $\varrho$ is well defined, invariant under $\mathfrak G_0$ and continuous on the compact set $\mathcal D_{\mathcal K_0}=\{U\in\mathfrak U(\theta):U^{(k)}\in\mathcal K_0\}$, and it vanishes exactly on the minimal orbits. For $V\in\mathcal K_0$,
\[
\mathscr N_r(V)=\{U\in\mathfrak B(V)\cap\mathfrak U(\theta):\ \varrho(U)<r\},
\]
a Borel set, and every open set of configurations that contains $\mathscr M(V)$ contains $\mathscr N_r(V)$ for $r$ small.
\end{enumerate}
\end{lemma}

\begin{proof}\hasdetail{lem:ol-chart-covariance}
(a) \emph{The hypotheses.} The regularity of Definition~\ref{def:regular} asks for local gauges, so it is invariant under real gauge transformations. Moreover $J_{U_0^u}=R(u)J_{U_0}$ (Lemma~\ref{lem:current}(b)), $W_j(U_0^u)=W_j^{u|_{\T^{(j)}}}$ for $j\le k$ (Lemma~\ref{lem:one-step-basic}(b)), and $V^v(c)\,W_k^v(c)^{-1}=v(c_-)\,V(c)W_k(c)^{-1}\,v(c_-)^{-1}$, whose logarithm has the same norm (Lemma~\ref{lem:log-estimates}(d)).

\emph{The averaging.} For a relative configuration $X$ at scale $j$, relative to $W_j=W_j(U_0)$, let $R(u)X$, $(R(u)X)(b)=R(u(b_-))X(b)$, be read relative to $W_j^u$, where $u$ also denotes its restriction to $\T^{(j)}$. Then $(R(u)X)W_j^u=(XW_j)^u$ bondwise. For $X$ at scale $j-1$, \eqref{eq:avg-contour-gauge} and Lemma~\ref{lem:one-step-basic}(b) give
\[
\begin{gathered}
(R(u)X)_{W^u}(y,x)=u(y)\,X_W(y,x)\,u(y)^{-1},\qquad\overline{R(u)X}=R(u)\,\overline X,\\
(R(u)X)^{(ugu^{-1})}=R(u)\big(X^{(g)}\big),
\end{gathered}
\]
with the pointwise conjugate $ugu^{-1}$; the last identity is $((XW)^u)^{ugu^{-1}}=((XW)^g)^u$. Hence $\tau_{R(u)X}=u\tau_Xu^{-1}$, the axial condition is preserved, $\mathcal E_j$, computed with $W^u$, satisfies $\mathcal E_j(uhu^{-1})=u\,\mathcal E_j(h)\,u^{-1}$, and $g_{R(u)X}=ug_Xu^{-1}$ by the uniqueness in Lemma~\ref{lem:axial}(a). So $\tilde{\mathfrak a}_j(R(u)X)=R(u)\,\tilde{\mathfrak a}_j(X)$, and by induction over $j$
\[
\mathcal Q_j\big(R(u)U_1;U_0^u\big)=R(u)\,\mathcal Q_j(U_1;U_0),\qquad\mathsf w_j\big(R(u)U_1;U_0^u\big)=u\,\mathsf w_j(U_1;U_0)\,u^{-1}.
\]
This is the argument of Proposition~\ref{prop:linearized-averages}(g), with $u$ in place of $r$, and it proves the nonlinear part of Proposition~\ref{prop:linearized-averages}(f). For $U_1=e^{\ii\eta A}$ with $|A|_\infty\le c_3$, all maps are defined at $U_1$ and at $R(u)U_1=e^{\ii\eta R(u)A}$ (Proposition~\ref{prop:Q-expansion}).

\emph{The gauge condition.} Directly from the definition, $D_{U_0^u}R(u)=R(u)D_{U_0}$ on scalar fields, and $Q'_k(U_0^u)R(u)=R(v)Q'_k(U_0)$ (Proposition~\ref{prop:linearized-averages}(f)). The maps $R(u)$ preserve the pairings. So every operator built from $D$, $D^T$, $Q'_k$ and $Q'^T_k$ by sums, products, transposes and inverses is covariant, as in the proof of Theorem~\ref{thm:variational}(h); in particular $R_{U_0^u}D^T_{U_0^u}R(u)=R(u)R_{U_0}D^T_{U_0}$.

\emph{The chart.} Let $A\in\mathcal S_V(U_0)$. Then $R(u)A$ is real, $|R(u)A|_\infty=|A|_\infty$, $R_{U_0^u}D^T_{U_0^u}R(u)A=0$, $e^{\ii\eta R(u)A}U_0^u=(e^{\ii\eta A}U_0)^u\in\mathfrak U(\theta)$, and $\mathcal Q(e^{\ii\eta R(u)A};U_0^u)=R(v)(VW_k^{-1})=V^v(W_k^v)^{-1}$. So $R(u)$ maps $\mathcal S_V(U_0)$ into $\mathcal S_{V^v}(U_0^u)$, and onto it, by the same argument for $u^{-1}$ at the reference $U_0^u$. With $\mathsf w=\mathsf w_k(e^{\ii\eta A};U_0)$,
\[
\mathsf U_{U_0^u}\big(R(u)A\big)=\big(R(u)e^{\ii\eta A}\big)^{(u\mathsf wu^{-1})}U_0^u=R(u)\big((e^{\ii\eta A})^{(\mathsf w)}\big)\,U_0^u=\big((e^{\ii\eta A})^{(\mathsf w)}U_0\big)^u=\mathsf U_{U_0}(A)^u .
\]

(b) Two elements of $\mathscr M(V)$ differ by some $u\in\mathfrak G_0$, for which $v=1$. By (a), the configurations $\mathsf U(A)$, $A\in\mathcal S_V$, $|A|_\infty<r$, for the reference $U_k(V)^u$ are the images under $u$ of those for $U_k(V)$, so they have the same $\mathfrak G_0$-orbits. Let now $u$ be any real gauge transformation on $\Tf$, and $v=u|_\T$. The curvature is covariant, so $V^v\in\mathfrak U_k(a_1)$; the minimal orbit of $V^v$ is $\mathscr M(V)^u$ (Theorem~\ref{thm:variational}(g), Theorem~\ref{thm:balaban}(a)); and $\mathfrak U(\theta')$ is invariant under real gauge transformations. So $(V^v,\theta)$ is admissible. With the reference $U_k(V)^u\in\mathscr M(V^v)$, (a) shows that the $\mathsf U(A')$, $A'\in\mathcal S_{V^v}$, $|A'|_\infty<r$, are the $\mathsf U(A)^u$, $A\in\mathcal S_V$, $|A|_\infty<r$. Since $(U^w)^u=(U^u)^{uwu^{-1}}$ and $u\mathfrak G_0u^{-1}=\mathfrak G_0$, the $\mathfrak G_0$-orbits of the $\mathsf U(A)^u$ are the images under $u$ of those of the $\mathsf U(A)$. Hence $\mathscr N_r(V^v)=\mathscr N_r(V)^u$.

(c) \emph{Definition.} Let $U\in\mathcal D_{\mathcal K_0}$, $V=U^{(k)}$ and $U_0\in\mathscr M(V)$. Since $\mathfrak U(\theta)\subset\mathfrak U_0(c_1)$, the $\mathfrak G_0$-orbit of $U$ lies in $\mathfrak B(V)\cap\mathfrak U(\theta)$. It has exactly one averaged axial representative relative to $U_0$ (Lemma~\ref{lem:axial}(c)), which is $\mathsf U_{U_0}(A)$ for exactly one $A\in\mathcal S_V(U_0)$ (Corollary~\ref{cor:parametrization}). For another reference $U_0^u$, $u\in\mathfrak G_0$, the configuration $\mathsf U_{U_0}(A)^u=\mathsf U_{U_0^u}(R(u)A)$ of (a) lies in the same orbit and is averaged axial relative to $U_0^u$, by Corollary~\ref{cor:parametrization}. So it is the representative relative to $U_0^u$, and $|R(u)A|_\infty=|A|_\infty$. Hence $\varrho$ is well defined, and it is constant on $\mathfrak G_0$-orbits. Since $\mathsf w_k(1)=1$, $\mathsf U_{U_0}(0)=U_0$, and $0\in\mathcal S_V(U_0)$, because $U_0\in\mathfrak U(\theta')$. So $\varrho(U)=0$ if and only if the representative is $U_0$, i.e.\ $U\in\mathscr M(V)$. The formula for $\mathscr N_r(V)$ restates its definition.

\emph{Continuity.} The class $\mathfrak U(\theta)$ is closed, hence compact, and $U\mapsto U^{(k)}$ is continuous on $\mathfrak U_0(c_1)$; so $\mathcal D_{\mathcal K_0}$ is compact. We use three facts.
\begin{itemize}
\item[(i)] The averages $U^{(j)}$, the averaged contour variables and the maps of Definition~\ref{def:adapted} are continuous functions of the configuration and of the reference, where they are defined: they are composed of group operations, $\ell$, $\exp$ and the group average $M$, which is holomorphic (Lemma~\ref{lem:group-average}(a)). So are $D_{U_0}$, $Q'_k(U_0)$ and $R_{U_0}=I-P_\cV$, with $P_\cV=T(T^TT)^{-1}T^T$, $T=G'Q'^T$ and $G'=(D^TD+aQ'^TQ')^{-1}$ (Section~\ref{sec:cov-scalar}); the matrices inverted are invertible at every reference considered. The averaged axial representative of the $\mathfrak G_0$-orbit of $U$ relative to $U_0$ is $U^g$, where $g|_\T=1$ and, scale by scale from $\T$ down to $\Tf$,
\[
g(x)=W_{i-1}(y,x)^{-1}\,g(y)\,U^{(i-1)}(y,x),\qquad x\in B_i(y),\ y\in\T^{(i)},
\]
with the averaged contour variables of $U^{(i-1)}$ and of $W_{i-1}=U_0^{(i-1)}$. This is the formula of Lemma~\ref{lem:axial}(a), applied as in the proof of Lemma~\ref{lem:axial}(b). So the representative depends continuously on $(U,U_0)\in\mathfrak U_0(c_1)^2$.
\item[(ii)] The union $\mathcal R$ of the minimal orbits $\mathscr M(V)$, $V\in\mathcal K_0$, is compact, and a limit of elements of $\mathscr M(V_n)$ with $V_n\to V$ lies in $\mathscr M(V)$. Indeed, let $U_n\in\mathscr M(V_n)$ with $V_n\in\mathcal K_0$. Pass to a subsequence with $V_n\to V\in\mathcal K_0$, and let $U_0\in\mathscr M(V)$. For $n$ large, $V_n=e^{\ii B_n}V$ bondwise, with real $B_n\to0$. Since $B_3a_1\le\frac12c_6$, Corollary~\ref{cor:analytic} applies with $\beta=\frac12c_6$: $\mathscr M(V_n)$ contains $\mathsf U_{U_0}(A_*(B_n))$, and $A_*(B_n)\to A_*(0)=0$. So $U_n=\mathsf U_{U_0}(A_*(B_n))^{u_n}$ with $u_n\in\mathfrak G_0$. Along a further subsequence $u_n\to u\in\mathfrak G_0$, and $U_n\to U_0^u\in\mathscr M(V)$, by the continuity of $\mathsf U_{U_0}$.
\item[(iii)] Let $\Sigma$ be the set of the pairs $(U_0,A)$ with $U_0\in\mathcal R$ and $A\in\mathcal S_{U_0^{(k)}}(U_0)$. By (i), the conditions that define $\mathcal S$ are closed in $(U_0,A)$, so $\Sigma$ is compact. The map $F(U_0,A)=(U_0,\mathsf U_{U_0}(A))$ is continuous, by (i), and injective, by Corollary~\ref{cor:parametrization}. So $F$ is a homeomorphism of $\Sigma$ onto its image.
\end{itemize}
Let $U_n\to U$ in $\mathcal D_{\mathcal K_0}$, and let $U_{0,n}\in\mathscr M(U_n^{(k)})$. By (ii), every subsequence has a further subsequence along which $U_{0,n}\to U_0\in\mathscr M(U^{(k)})$. Along it, the representatives $\mathsf U_{U_{0,n}}(A_n)$ of the orbits of the $U_n$ converge to the representative $\mathsf U_{U_0}(A)$ of the orbit of $U$, by (i), and $A_n\to A$, by (iii). So $\varrho(U_n)\to\varrho(U)$ along it. Hence $\varrho$ is continuous.

\emph{Neighbourhoods.} Let $V\in\mathcal K_0$. The set $\mathscr N_r(V)$ is relatively open in the compact set $\mathfrak B(V)\cap\mathfrak U(\theta)=\{U\in\mathfrak U(\theta):U^{(k)}=V\}$, hence a Borel set. Let $\mathcal V$ be open with $\mathscr M(V)\subset\mathcal V$. The sets $K_r=\{U\in\mathfrak B(V)\cap\mathfrak U(\theta):\varrho(U)\le r\}$ are compact, they decrease as $r\downarrow0$, and their intersection is $\mathscr M(V)$. So $K_r\setminus\mathcal V$ is empty for $r$ small, and then $\mathscr N_r(V)\subset K_r\subset\mathcal V$.
\end{proof}

\begin{lemma}[Fibre mass]\label{lem:ol-fibre-mass}
Let $(V,\theta_0)$ be admissible. Then $U\mapsto U^{(k)}$ is a submersion at every point of $\mathfrak B(V)\cap\mathfrak U(\theta_0)$, $\mathfrak B(V)$ is an analytic submanifold near this compact set, and the measure $\prod_bdU(b)\prod_c\delta(U^{(k)}(c)V(c)^{-1})$ gives it a finite mass. If $\mathcal K_0$ is a compact set of real configurations on $\T$ such that $(V,\theta_0)$ is admissible for every $V\in\mathcal K_0$, these masses are bounded uniformly for $V\in\mathcal K_0$.
\end{lemma}

\begin{proof}\hasdetail{lem:ol-fibre-mass}
The set $\mathfrak B(V)\cap\mathfrak U(\theta_0)$ is compact, since $U\mapsto U^{(k)}$ is continuous on $\mathfrak U_0(c_1)$. Since $(V,\theta_0)$ is admissible, the setting applies with $\theta_0$ in place of $\theta$ and a reference $U_0\in\mathscr M(V)$, and Corollary~\ref{cor:parametrization} applies with $\eps=B_3a_1$ and $\beta=0$. By Theorem~\ref{thm:chart} every $A\in\mathcal S_V$ satisfies $|A|_{C^1}\le B_1(\eps+\theta_0)=\frac12\rho_1$, so $\mathcal S_V$ is a compact subset of the manifold $\mathcal M_V$ of Lemma~\ref{lem:reduction}(a), near which the map $A\mapsto\ell(\mathcal Q(e^{\ii\eta A}))=\ell(\mathsf U(A)^{(k)}W_k^{-1})$ (for every $A$ with $|A|_\infty\le c_3$, by Theorem~\ref{thm:correspondence}(a), as in the proof of Corollary~\ref{cor:parametrization}) on the slice $\{RD^TA=0\}$ is a submersion (proof of Lemma~\ref{lem:reduction}(a)). Hence $U\mapsto U^{(k)}$ is a submersion at every $\mathsf U(A)$, $A\in\mathcal S_V$, and, since $(U^u)^{(k)}=U^{(k)}$ for $u\in\mathfrak G_0$, at every point of the compact set $\mathfrak B(V)\cap\mathfrak U(\theta_0)$ (Corollary~\ref{cor:parametrization}). So $\mathfrak B(V)$ is an analytic submanifold near this set, the Haar $\delta$-measure has a continuous density on it there, and the mass is finite.

For the uniform bound, the set $\mathcal D=\{U\in\mathfrak U(\theta_0):U^{(k)}\in\mathcal K_0\}$ is compact, as in Lemma~\ref{lem:ol-chart-covariance}(c), and $U\mapsto U^{(k)}$ is a submersion at each of its points, by the first part. Near each point of $\mathcal D$ it is a coordinate projection in suitable analytic coordinates, in which the measures $\prod_bdU(b)\prod_c\delta(U^{(k)}(c)V(c)^{-1})$ on the fibres have bounded densities; finitely many such coordinate neighbourhoods cover $\mathcal D$. So the masses of the sets $\mathfrak B(V)\cap\mathfrak U(\theta_0)$, $V\in\mathcal K_0$, are bounded by some $m$.
\end{proof}

\begin{proposition}[The one-loop term]\label{prop:ol-laplace}
Let $\theta\le\theta_0$. There are $r_0>0$, constants $c_k(g)>0$ independent of $V$, $r$ and $\theta$, and a function $\Gamma^{(1)}_k(V)$ independent of $r$, $g$ and $\theta$, such that for $0<r\le r_0$ and every real $V\in\mathfrak U_k(a_1)$ for which $(V,\theta)$ is admissible
\[
I_k(V;g,r)=c_k(g)\exp\Big(-\frac1{g^2}S\big(U_k(V)\big)+\Gamma^{(1)}_k(V)+o(1)\Big)\qquad(g\to0),
\]
locally uniformly in $V$ on the set of these $V$: every point of this set has a neighbourhood in it on which the $o(1)$ is uniform. $\Gamma^{(1)}_k$ is real analytic in $V$ and invariant under gauge transformations of $\T$.
\end{proposition}

\begin{proof}\hasdetail{prop:ol-laplace}
\emph{Plan.} Modulo $\mathfrak G_0$, the integral $I_k$ runs over a neighbourhood of the minimal orbit in the constraint set, which the chart $\mathcal S_V$ parametrizes. On $\mathcal S_V$ the action $\Phi$ has the unique minimum $A=0$, with Hessian $q_2$ on the tangent space; Laplace's method there gives $\Gamma^{(1)}_k$. A reference fixed for all $V$ near a given $V^\circ$ gives the uniformity in $V$.

\emph{The constraint set.} Near $\mathfrak B(V)$, the differential of $U\mapsto U^{(k)}$ in the exponential coordinates is $\bQ_k$ plus $D_V$ composed with the linearization of $\mathsf w_k$ at the points of $\T$ (Lemma~\ref{lem:Q-covariance}(a)). It is onto: by the chain rule applied to $A\mapsto\mathsf U(A)^{(k)}=\mathcal Q(e^{\ii\eta A})W_k$ (Theorem~\ref{thm:correspondence}(a)), whose differential at $0$ on the slice is $\bQ_k$, with $\bQ_k\bH=I$. So $\mathfrak B(V)$ is a real analytic submanifold near the minimal orbit, and, by Lemma~\ref{lem:ol-fibre-mass}, near every point of $\mathfrak B(V)\cap\mathfrak U(\theta)$. The Haar $\delta$-function disintegrates the Haar measure into a positive analytic measure on it.

\emph{The quotient.} The group $\mathfrak G_0$ is compact and acts freely, so the quotient by $\mathfrak G_0$ of a neighbourhood of $\mathfrak B(V)\cap\mathfrak U(\theta)$ in $\mathfrak B(V)$ is a real analytic manifold, and $\mathfrak B(V)\cap\mathfrak U(\theta)$ modulo $\mathfrak G_0$ is a compact subset of it. By Corollary~\ref{cor:parametrization}, $A\mapsto[\mathsf U(A)]$ is a bijection from $\mathcal S_V$ onto this compact set. It is a chart of the quotient near $A=0$; see the paragraph ``The chart at $0$'' below. Since $U_0\in\mathfrak U(\theta')$ with $\theta'<\theta$, and curvature and current are continuous, $\mathcal S_V$ contains a neighbourhood of $0$ in the manifold of the real $A$ with $RD^TA=0$ and $\mathcal Q(e^{\ii\eta A})=1$ (Lemma~\ref{lem:reduction}(a)). In this chart the integrand is $e^{-\Phi(A)/g^2}$, and near $A=0$ the measure has a positive analytic density $\varsigma_V(A)$ with respect to the Euclidean volume on $\mathcal S_V$ (paragraph ``The chart at $0$''). The density includes the Jacobian of $d\mathsf U(0)$ and the volume of the orbits. Laplace's method uses $\varsigma_V$ only near $A=0$; away from $A=0$ only the finiteness of the mass is used (Lemma~\ref{lem:ol-fibre-mass}).

\emph{The minimum.} By Proposition~\ref{prop:minimality}, $\Phi$ has on $\mathcal S_V$ the unique minimum $A=0$. Its Hessian on the tangent space $\cX=\{h:\bQ_kh=0,\ RD^Th=0\}$ is $q_2$ (Proposition~\ref{prop:hessian}): $0$ is critical on the constraint set, and the curvature of the constraint gives the multiplier term. It is coercive. Indeed, the reference $U_0=U_k(V)$ is critical, with $\beta=0$ and $A_*=0$, and it is $(M,\eps)$-regular with $|J_{U_0}|\le\eps$ for $\eps=B_3a_1$ (Theorem~\ref{thm:balaban}(b) with $\eps_1=a_1$), where $\eps\le\frac12c_6\le\frac12c_7$. On all real fine fields $h$, $q_2(h,h)=\|Dh\|^2+d^2\mathcal N(0)[h,h]-\ip{\omega_0}{d^2\Cav_k(0)[h,h]}$, with $|d^2\mathcal N(0)[h,h]|\le C\eps(\|h\|^2+\|\nabla h\|^2)$, $|\ip{\omega_0}{d^2\Cav_k(0)[h,h]}|\le C\eps\|h\|^2$ and the multiplier $\omega_0$ of Proposition~\ref{prop:euler-lagrange} (proof of Lemma~\ref{lem:Lstar-coercive}). By Lemma~\ref{lem:Lstar-coercive}, $q_2(h,h)\ge\frac12c_*(\|h\|^2+\|Dh\|^2)$ for $h\in\cX$. Hence $\Phi\ge\Phi(0)+cr^2$ on $\{|A|_\infty=r\}\cap\mathcal S_V$ for $r\le r_1$, with $r_1,c>0$ depending on $V$. Since $\mathcal S_V$ is compact and $0$ is the unique minimum, $\Phi-\Phi(0)\ge m_1>0$ on $\mathcal S_V\cap\{|A|_\infty\ge r_1\}$, so $\Phi\ge\Phi(0)+c'|A|_\infty^2$ on all of $\mathcal S_V$, with $c'=\min(c,m_1/r_u^2)$. No smallness of $r$ is used below; we take $r_0=r_u$.

\emph{The chart at $0$.} Let $h\in\cX$, and let $A(t)$ be a curve in the manifold $\{RD^TA=0,\ \mathcal Q(e^{\ii\eta A})=1\}$ with $A(0)=0$ and $A'(0)=h$ (Lemma~\ref{lem:reduction}(a)). Then $\gamma(t)=\mathsf U(A(t))$ is a curve in $\mathfrak B(V)$ through $U_0$, $\gamma'(0)=d\mathsf U(0)h$, and $S(\gamma(t))=\Phi(A(t))$, whose second derivative at $t=0$ is $q_2(h,h)$, since $0$ is critical on the constraint set with the multiplier $\omega_0$. Since $U_0$ minimizes $S$ on $\mathfrak B(V)\cap\mathfrak U(\theta)$ (Theorem~\ref{thm:balaban}(a)), it is a critical point of $S|_{\mathfrak B(V)}$, and the second derivative of $S$ along a curve in $\mathfrak B(V)$ through $U_0$ is the Hessian of $S|_{\mathfrak B(V)}$ at $U_0$, evaluated at the velocity. Since $S$ is constant on the $\mathfrak G_0$-orbit of $U_0$, this Hessian vanishes on the tangent vectors of the orbit. So if $d\mathsf U(0)h$ is tangent to the orbit, then $q_2(h,h)=0$ and $h=0$. Hence $(h,\xi)\mapsto d\mathsf U(0)h+D_{U_0}\xi$, $\xi|_\T=0$, is injective, since $D_{U_0}\xi=0$ and $\xi|_\T=0$ give $\xi=0$. The map $(A,u)\mapsto\mathsf U(A)^u$, on a neighbourhood of $0$ in the constraint manifold times $\mathfrak G_0$, is therefore an injective immersion: it is an immersion at the points $(0,u)$, by the $\mathfrak G_0$-equivariance, hence near them, since $\mathfrak G_0$ is compact; and it is injective by Corollary~\ref{cor:parametrization}, since $\mathfrak G_0$ acts freely. Its image contains a neighbourhood of $\mathscr M(V)$ in $\mathfrak B(V)$ (Lemma~\ref{lem:ol-chart-covariance}(c) and the uniqueness in the implicit function theorem). So the dimensions agree, since otherwise the image would have measure zero, and the map is an analytic diffeomorphism onto an open set. In particular $A\mapsto[\mathsf U(A)]$ is a chart near $0$, and the density $\varsigma_V$ is positive.

\emph{Laplace's method} on the finite-dimensional manifold $\mathcal S_V$ gives the expansion, with
\[
\Gamma^{(1)}_k(V)=-\tfrac12\log\det\big(q_2|_{\cX}\big)+\log\varsigma_V(0)
\]
up to a constant, and $c_k(g)$ is $(2\pi g^2)^{N/2}$, $N=\dim\cX$, times a constant. Both terms are local at the minimum, so they do not depend on $r$, nor on the choice of $\theta$, and they are analytic in $V$ (paragraph ``Uniformity in $V$'', (vi)). For the gauge invariance let $u$ be a real gauge transformation on $\Tf$ and $v=u|_\T$. By Lemma~\ref{lem:ol-chart-covariance}(b), $(V^v,\theta)$ is admissible and $\mathscr N_r(V^v)=\mathscr N_r(V)^u$, and $V^v\in\mathfrak U_k(a_1)$. The substitution $U\mapsto U^u$ preserves $\prod_bdU(b)$ and $S$, and $\delta((U^u)^{(k)},V^v)=\delta(U^{(k)},V)$, by Lemma~\ref{lem:one-step-basic}(b) and since $\delta$ is a class function. So $I_k(V^v;g,r)=I_k(V;g,r)$. Since $S(U_k(V^v))=S(U_k(V))$ and $c_k(g)$ does not depend on $V$, $\Gamma^{(1)}_k(V^v)=\Gamma^{(1)}_k(V)$.

\emph{Uniformity in $V$.} So far $V$ was fixed. Let $V^\circ$ be a point of the set of the statement, and fix $U_0\in\mathscr M(V^\circ)$. For $V$ near $V^\circ$ write $V(c)=e^{\ii B(c)}V^\circ(c)$; then $B=B(V)$ is real, small and analytic in $V$. For all these $V$ we use the fixed reference $U_0$, with its operators $R$, $D$, $\mathcal Q$, $\bQ_k$, its minimizer map $\bH$, the space $\cX$ and $\Phi(A)=S(e^{\ii\eta A}U_0)$. By Corollary~\ref{cor:parametrization}, with $\beta=|B|_\infty$, $A\mapsto\mathsf U_{U_0}(A)$ is a bijection from the set $\mathcal S_V(U_0)$ of the real $A$ with $|A|_\infty\le r_u$, $RD^TA=0$, $e^{\ii\eta A}U_0\in\mathfrak U(\theta)$ and $\mathcal Q(e^{\ii\eta A})=e^{\ii B}$ onto the averaged axial representatives of the $\mathfrak G_0$-orbits in $\mathfrak B(V)\cap\mathfrak U(\theta)$, and $S(\mathsf U_{U_0}(A))=\Phi(A)$. By Corollary~\ref{cor:analytic}, the minimal orbit of $V$ is the $\mathfrak G_0$-orbit of $\mathsf U_{U_0}(A_*(B))$, with $A_*$ analytic, $A_*(0)=0$, $RD^TA_*(B)=0$ and $\mathcal Q(e^{\ii\eta A_*(B)})=e^{\ii B}$.
\begin{itemize}
\item[(i)] \emph{Coordinates.} Fix a Euclidean norm $|\cdot|$ on $\cX$. On the slice $\{RD^TA=0\}$ the map $f(A)=\ell(\mathcal Q(e^{\ii\eta A}))$ is analytic, with $f(0)=0$ and $df(0)=\bQ_k$ (Proposition~\ref{prop:Q-expansion}); $\bH$ maps into the slice, and $\bQ_k\bH=I$ (proof of Lemma~\ref{lem:reduction}(a)). By the analytic implicit function theorem, for $(x,B)$ near $(0,0)$ the equation $f(x+\bH y)=B$ has a unique small solution $y=y(x,B)$, analytic in $(x,B)$. So $A(x,B)=x+\bH y(x,B)$ parametrizes $\mathcal S_V(U_0)$ near $0$ by a neighbourhood of $0$ in the fixed space $\cX$, and $x=(I-\bH\bQ_k)A(x,B)$. Put $x_*(B)=(I-\bH\bQ_k)A_*(B)\in\cX$; by the uniqueness, $A(x_*(B),B)=A_*(B)$.
\item[(ii)] \emph{The phase.} $F(x,B)=\Phi(A(x,B))$ is analytic. Since $A_*(B)$ is a critical point of $\Phi$ on the constraint set (Lemma~\ref{lem:reduction}(b)), $x_*(B)$ is a critical point of $F(\cdot,B)$, and $F(x_*(B),B)=S(U_k(V))$. Its Hessian $H(B)=\partial_x^2F(x_*(B),B)$ is continuous in $B$, and $H(0)$ is $q_2$ on $\cX$, which is positive definite (paragraph ``The minimum''). Since the third derivatives of $F$ are bounded on a compact neighbourhood of $(0,0)$, there are $\delta,c_H>0$ such that, for $B$ small, $F(x,B)-F(x_*(B),B)\ge c_H|x-x_*(B)|^2$ for $|x-x_*(B)|<\delta$.
\item[(iii)] \emph{The density.} Let $\Upsilon(x,B,u)=\mathsf U_{U_0}(A(x,B))^u$, for $(x,B)$ near $(0,0)$ and $u\in\mathfrak G_0$. It is analytic, and $\Upsilon(x,B,u)^{(k)}=e^{\ii B}V^\circ$, by Corollary~\ref{cor:parametrization} and $(U^u)^{(k)}=U^{(k)}$. At $B=0$, $(x,u)\mapsto\Upsilon(x,0,u)$ is a local diffeomorphism into $\mathfrak B(V^\circ)$ near $\mathscr M(V^\circ)$, by the paragraph ``The chart at $0$'' at $V^\circ$. If $d\Upsilon$ annihilates a tangent vector at a point $(0,0,u)$, then the $B$-component of the vector vanishes, since $B\mapsto e^{\ii B}V^\circ$ is a local diffeomorphism, and then the vector vanishes. Since $U\mapsto U^{(k)}$ is a submersion near $\mathscr M(V^\circ)$, the space of the $(x,B,u)$ has the dimension of the space of configurations, so $d\Upsilon$ is invertible at these points, and, by the compactness of $\mathfrak G_0$, on $\mathcal P\times\mathfrak G_0$ for a neighbourhood $\mathcal P$ of $(0,0)$. Moreover $\Upsilon$ is injective there: equal values have equal averages, hence equal $B$; then they are gauge transforms of averaged axial representatives of the same $\mathfrak G_0$-orbit, which are equal (Lemma~\ref{lem:axial}(c)), so $A(x,B)$ and $x$ are equal (Corollary~\ref{cor:parametrization}), and so is $u$, since the action is free. Hence $\Upsilon$ is an analytic diffeomorphism of $\mathcal P\times\mathfrak G_0$ onto an open set of configurations. Pulled back by $\Upsilon$, the Haar measure has a positive analytic density $\mathsf h(x,B,u)$ with respect to $dx\,dB\,du$, with the Haar measure $du$ of $\mathfrak G_0$. Since $\Upsilon^{(k)}=e^{\ii B}V^\circ$, the measure $\prod_bdU(b)\prod_c\delta(U^{(k)}(c)V(c)^{-1})$ on the image is the image of $j(B)\mathsf h(x,B,u)\,dx\,du$ at $B=B(V)$, with $j(B)=\prod_c\sigma(B(c))^{-1}$, where $\sigma$ is the density of the normalized Haar measure of $G$ in exponential coordinates (Section~\ref{subsec:exact-formula}). So the integral of a $\mathfrak G_0$-invariant function $\mathcal F$ over $\Upsilon(\mathcal P\times\mathfrak G_0)\cap\mathfrak B(V)$ is $\int\mathcal F(\mathsf U_{U_0}(A(x,B)))\,\varsigma(x,B)\,dx$, with $\varsigma(x,B)=j(B)\int_{\mathfrak G_0}\mathsf h(x,B,u)\,du$, positive and analytic.
\item[(iv)] \emph{Laplace's method with parameters.} For $V$ near $V^\circ$ and $B=B(V)$ let $\mathcal P_\delta(V)=\Upsilon\big(\{x:|x-x_*(B)|<\delta\}\times\{B\}\times\mathfrak G_0\big)\subset\mathfrak B(V)$. By (iii), the integral of $e^{-S/g^2}$ over $\mathcal P_\delta(V)$, for the measure of $I_k$, is $\int_{|x-x_*(B)|<\delta}e^{-F(x,B)/g^2}\varsigma(x,B)\,dx=g^Ne^{-F(x_*(B),B)/g^2}\int\phi_g(y,B)\,dy$, where $\phi_g(y,B)=e^{-(F(x_*+gy,B)-F(x_*,B))/g^2}\varsigma(x_*+gy,B)$, with $x_*=x_*(B)$, for $g|y|<\delta$, and $\phi_g(y,B)=0$ otherwise. By (ii), $0\le\phi_g(y,B)\le Ce^{-c_H|y|^2}$. Since $\partial_xF(x_*(B),B)=0$ and $F$ and $\varsigma$ have bounded derivatives on a compact neighbourhood of $(0,0)$, $\phi_g(y,B)\to e^{-\frac12H(B)[y,y]}\varsigma(x_*(B),B)$ as $g\to0$, uniformly for $y$ in compact sets and $B$ small. With the uniform Gaussian bound on the tails, $\int\phi_g(y,B)\,dy\to(2\pi)^{N/2}\det H(B)^{-1/2}\varsigma(x_*(B),B)$, uniformly for $B$ small. So the integral over $\mathcal P_\delta(V)$ is
\[
(2\pi g^2)^{N/2}e^{-S(U_k(V))/g^2}\det H(B)^{-1/2}\varsigma(x_*(B),B)\big(1+o(1)\big),
\]
uniformly for $V$ near $V^\circ$.
\item[(v)] \emph{The rest of the integral.} Let $b_0>0$ be small, and let $\mathcal K_{V^\circ}$ be the set of the $V=e^{\ii B}V^\circ\in\mathfrak U_k(a_1)$ with $|B|_\infty\le b_0$. It is compact, and it is a neighbourhood of $V^\circ$ in the set of the statement: for $b_0$ small, $(V,\theta)$ is admissible for every $V\in\mathcal K_{V^\circ}$, since its minimal orbit consists of gauge transforms of $e^{\ii\eta A_*(B)}U_0$, whose curvature and current are close to those of $U_0\in\mathfrak U(\theta')$. Let $\mathcal D=\{U\in\mathfrak U(\theta):U^{(k)}\in\mathcal K_{V^\circ}\}$, a compact set. The function $U\mapsto S(U)-S(U_k(U^{(k)}))$ is continuous on $\mathcal D$, since $S(U_k(e^{\ii B}V^\circ))=\Phi(A_*(B))$; it is nonnegative, and it vanishes exactly on the minimal orbits (Proposition~\ref{prop:minimality}). The set $\mathcal O_\delta=\Upsilon\big(\{(x,B):|x-x_*(B)|<\delta,\ |B|_\infty<2b_0\}\times\mathfrak G_0\big)$ is open and contains these orbits, and a point of $\mathfrak B(V)$ in $\mathcal O_\delta$ lies in $\mathcal P_\delta(V)$, since $\Upsilon^{(k)}=e^{\ii B}V^\circ$. The function $\varrho$ of Lemma~\ref{lem:ol-chart-covariance}(c), with $\mathcal K_{V^\circ}$ in place of $\mathcal K_0$, is continuous on $\mathcal D$ and vanishes exactly on these orbits. So there is $c>0$ such that $S(U)\ge S(U_k(U^{(k)}))+2c$ on the compact sets $\mathcal D\setminus\mathcal O_\delta$ and $\{U\in\mathcal D:\varrho(U)\ge r\}$. Moreover $(V,\theta_0)$ is admissible for every $V\in\mathcal K_{V^\circ}$, since $\theta\le\theta_0$. So, by Lemma~\ref{lem:ol-fibre-mass}, the masses of the sets $\mathfrak B(V)\cap\mathfrak U(\theta)\subset\mathfrak B(V)\cap\mathfrak U(\theta_0)$, $V\in\mathcal K_{V^\circ}$, for the measure of $I_k$, are bounded by some $m$. For $\delta$ and $b_0$ small, $\mathcal P_\delta(V)\subset\mathfrak B(V)\cap\mathfrak U(\theta)$, as for the minimal orbits: its points are gauge transforms of $e^{\ii\eta A(x,B)}U_0$ with $A(x,B)$ small, so on the fixed torus their curvature and current are within $o(1)$ of those of $U_0\in\mathfrak U(\theta')$, $\theta'<\theta$. Now $I_k(V;g,r)$ is the integral over $\mathscr N_r(V)=\{U\in\mathfrak B(V)\cap\mathfrak U(\theta):\varrho(U)<r\}$ (Lemma~\ref{lem:ol-chart-covariance}(c)). It differs from the integral over $\mathcal P_\delta(V)$ by at most the integrals over $\mathscr N_r(V)\setminus\mathcal P_\delta(V)\subset\mathcal D\setminus\mathcal O_\delta$ and over $\mathcal P_\delta(V)\setminus\mathscr N_r(V)\subset\{U\in\mathcal D:\varrho(U)\ge r\}$, hence by at most $me^{-(S(U_k(V))+2c)/g^2}$. By (iv) this is a relative error $O(g^{-N}e^{-2c/g^2})$, uniformly for $V\in\mathcal K_{V^\circ}$.
\item[(vi)] \emph{Conclusion.} By (iv) and (v), uniformly for $V\in\mathcal K_{V^\circ}$,
\[
I_k(V;g,r)=(2\pi g^2)^{N/2}e^{-S(U_k(V))/g^2}\det H(B)^{-1/2}\varsigma(x_*(B),B)\big(1+o(1)\big).
\]
At each single $V\in\mathcal K_{V^\circ}$ the expansion of the paragraphs ``The constraint set''--``Laplace's method'' holds as well (it is (iv)--(v) with $V^\circ=V$, $B=0$), and the terms of order $g^0$ of the two agree. Since $c_k(g)$ is $(2\pi g^2)^{N/2}$ times a constant, $\Gamma^{(1)}_k(V)=-\frac12\log\det H(B(V))+\log\varsigma(x_*(B(V)),B(V))$ on $\mathcal K_{V^\circ}$, up to a constant, and the $o(1)$ of the statement is uniform on $\mathcal K_{V^\circ}$.
\end{itemize}
\end{proof}

\begin{lemma}[Choice of the neighbourhood]\label{lem:ol-neighbourhood}
Let $V\in\mathfrak U_k(a_1)$ and $\theta\le\theta_0$ be such that $(V,\theta)$ is admissible (for instance $V$ small field and $\theta$ as above), let $r\le r_0$, and let $\mathscr N$ be a Borel set with $\mathscr N_r(V)\subset\mathscr N\subset\mathfrak B(V)\cap\mathfrak U(\theta_0)$. Then $I_k(V;g,r)$ and the same integral over $\mathscr N$ differ by a factor $1+O(e^{-c/g^2})$, with $c>0$. In particular their terms of order $g^0$ agree.
\end{lemma}

No invariance of $\mathscr N$ under $\mathfrak G_0$ is required.

\begin{proof}\hasdetail{lem:ol-neighbourhood}
Put $S_*=S(U_k(V))$. The set $\mathfrak B(V)\cap\mathfrak U(\theta_0)$ is compact, since $U\mapsto U^{(k)}$ is continuous on $\mathfrak U_0(c_1)$. Its part outside $\mathscr N_r(V)$ lies in $K_1\cup K_2$. Here $K_1$ is the set of its elements whose curvature or current has a value of norm at least $\theta$, and $K_2=\{U\in\mathfrak B(V)\cap\mathfrak U(\theta):\varrho(U)\ge r\}$, with $\varrho$ of Lemma~\ref{lem:ol-chart-covariance}(c) for $\mathcal K_0=\{V\}$. Both sets are compact and invariant under $\mathfrak G_0$, since curvature and current are gauge covariant, and by that lemma. They do not meet $\mathscr M(V)$, since $\mathscr M(V)\subset\mathfrak U(\theta')$ with $\theta'<\theta$, and $\varrho=0$ on $\mathscr M(V)$. By Proposition~\ref{prop:minimality}, $S>S_*$ on $\mathfrak B(V)\cap\mathfrak U(\theta_0)$ outside $\mathscr M(V)$. So $S\ge S_*+2c$ on $K_1\cup K_2$, for some $c>0$, and the integral over $\mathscr N\setminus\mathscr N_r(V)$ is at most $e^{-(S_*+2c)/g^2}m$, where $m$ is the mass of $K_1\cup K_2$ for the measure $\prod_bdU(b)\prod_c\delta(U^{(k)}(c)V(c)^{-1})$. This mass is finite, by Lemma~\ref{lem:ol-fibre-mass}, since the pair $(V,\theta_0)$ is admissible and $K_1\cup K_2\subset\mathfrak B(V)\cap\mathfrak U(\theta_0)$. On the other hand, by Proposition~\ref{prop:ol-laplace}, $I_k(V;g,r)\ge\frac12c_k(g)\exp(-S_*/g^2+\Gamma^{(1)}_k(V))$ for $g$ small, and $c_k(g)$ is a constant times a power of $g$, by Laplace's method. So the two integrals differ by a factor $1+O(e^{-c/g^2})$.
\end{proof}

\paragraph{The one-loop terms of the steps.} For the step from scale $j$ to scale $j+1$, in the units of step $j$, Theorem~\ref{thm:step} gives for the gauge-fixed step $T^{\rm gf}(\chi^{\rm w}\rho)$ the prefactor $\mathcal N_j\exp(\Gamma^{\rm step}_j(U'))$, where
\[
\Gamma^{\rm step}_j(U')=-\tfrac12\log\det\Delta_j+\sum_c\log\mathfrak j_c ,\qquad U'=U_{j+1}(V'),
\]
with $\Delta_j$ the fluctuation operator of the step on the variables $\mathfrak F$ (Lemma~\ref{lem:step-quadratic}, Proposition~\ref{prop:det-densities}) and $\mathfrak j_c$ the Jacobians of Theorem~\ref{thm:step}(b). The operator $\Delta_j$ contains the quadratic part $\|\mathfrak G_W\cdot\|^2$ of the gauge fixing. Both terms are functions of the background $U'$, a configuration on the finest lattice, and $\mathcal N_j$ does not depend on it. Up to a constant, $\Gamma^{\rm step}_j$ does not depend on the gauge fixing. Remark~\ref{rem:one-loop-intrinsic} sketches this, and the sketch is not used below; Proposition~\ref{prop:ol-composition} below proves the composite version of this statement. The Jacobians are not negligible: $\log\mathfrak j_c$ is quadratic in the curvature of scale $j+1$ (Proposition~\ref{prop:jacobians}(c)), so it contributes to the marginal tensor below. For $j=0$ the base is $W=U'$ and the background map is trivial. The formula of Theorem~\ref{thm:step} then holds with the fluctuation operator $\Delta_0$ of Proposition~\ref{prop:first-step}(b): this is Lemma~\ref{lem:first-step-formula}, the \emph{formula for the first step}. It is used by Proposition~\ref{prop:ol-composition} through its case $k=1$, and hence by Corollary~\ref{cor:ol-cumulative}; its part~(c) is used by Lemma~\ref{lem:ol-symmetry}(b), and hence~(c), for the step from scale $0$. Lemma~\ref{lem:ol-one-step}, which concerns the steps from the scales $j\ge1$, does not use it. In the flow of Part~X the same function appears as the one-loop part of the new generation: $\Gamma^{\rm step}_j(U')-\Gamma^{\rm step}_j(1)=\sum_X\mathfrak e(X,U')+\sum_c\log\mathfrak j_c(U')$ (Proposition~\ref{prop:one-loop-polymer}(a); Proposition~\ref{prop:first-step}(c) for $j=0$) is the sum of the one-loop part $E'_{\rm ol}$ of the raw new generation of the step (Definition~\ref{def:raw-new}, Section~\ref{sec:ol-uv}). It involves neither the coupling nor the old generations. Its coefficient, read off on the torus, is the one-loop part $\beta^{\rm ol}_{j+1}$ of $\beta_{j+1}$, and Proposition~\ref{prop:flow-F} compares it with the coefficient of the marginal tensor below.

\begin{lemma}[The one-step relation]\label{lem:ol-one-step}
Let $k\ge2$, let $V$ be small field, and put $V^0=U_k(V)^{(k-1)}$, a configuration of scale $k-1$, read in the units of step $k-1$. Then $V^0$ is small field at scale $k-1$, and, up to a constant independent of $V$,
\[
\Gamma^{(1)}_k(V)=\Gamma^{(1)}_{k-1}\big(V^0\big)+\Gamma^{\rm step}_{k-1}\big(U_k(V)\big).
\]
\end{lemma}

\begin{proof}\hasdetail{lem:ol-one-step}
\emph{Plan.} Integrate out the configuration of scale $k-1$ (``Factorization''). The inner integral is $I_{k-1}$, which Proposition~\ref{prop:ol-laplace} evaluates at scale $k-1$. In the outer integral, replace the region by the fluctuation cutoff $\chi^{\rm w}$, by a two-radius estimate (the terms $R_1$, $R_2$). Then Theorem~\ref{thm:step} applies to the outer integral, and the result is compared with Proposition~\ref{prop:ol-laplace} at scale $k$.

\emph{Notation.} We use the notation of Sections~\ref{sec:fluctuations} and~\ref{sec:covariance} for the step from scale $k-1$ to scale $k$, in the units of step $k-1$; in particular, in this proof $\T$ and $\T'$ denote the lattices of scale $k-1$ and $k$, and $B'(y)$, $y\in\T'$, the blocks of the step. So $V'=V$, $U'=U_k(V)$, the base is $W=V^0=U'^{(k-1)}$, and $\mathcal G'$ is the group of real gauge transformations of the lattice of scale $k-1$ that are $1$ on the lattice of scale $k$. Let $c_L$ be the constant of Theorem~\ref{thm:step-coercive}, and $r_2$, $C_L$ those of Theorem~\ref{thm:step}(d). We take $w\le\min(r_2,\,c_L/(16C_L))$, and smaller where stated. We also take $8w\le a'_1$ and $16B_3w\le L^{-3}\theta_0$ (both implied by $w\le r_2$ if $c\le\frac1{16}$ in Lemma~\ref{lem:step-quadratic}); then the same holds for $w'<w$. Put $C_1=128C_M$, with $C_M$ from Lemma~\ref{lem:group-average}, and $w'=w/(C_1L)$.

\emph{The region.} Let $\mathscr A$ be the set of configurations $V_{k-1}$ on $\T$ such that, for every bond $b$, $V_{k-1}(b)=e^{\ii Z(b)}W(b)$ with $Z(b)\in\mathfrak g$ and $|Z(b)|<w'$, and let $\mathscr O=\{V_{k-1}^u:\ V_{k-1}\in\mathscr A,\ u\in\mathcal G'\}$ be its $\mathcal G'$-saturation. By \eqref{eq:chi-w-support}, with $w'$ in place of $w$, $\mathscr A\subset\mathfrak M$. Since $\mathfrak M$ is invariant under real gauge transformations, $\mathscr O\subset\mathfrak M$, and $\mathscr O$ is $\mathcal G'$-invariant. A gauge transformation acts on each bond variable by a left and a right translation, which preserve the bi-invariant distance on $G$. So every point of $\mathscr O$ lies within $w'$, bondwise, of the orbit $\{(V^0)^u:u\in\mathcal G'\}$. Let $\eps'_1\le a_{\rm step}$ bound the curvature of $V$ in the units of step $k$, as in Theorem~\ref{thm:step}. By Lemma~\ref{lem:composition}(a) and $4(C_*+1)L^{-2}a_{\rm step}\le a'_1$ (proof of Theorem~\ref{thm:step}), $W\in\mathfrak U_{k-1}(\frac12a'_1)$, with $a'_1$ of Lemma~\ref{lem:composition}(c). Let $\mathcal K_0$ be the $\mathcal G'$-saturation of the set of the $e^{\ii Z}W$ with $|Z(b)|\le w'$ on all bonds. It is compact and contains $\mathscr O$, and, by the argument before \eqref{eq:chi-w-support}, $\mathcal K_0\subset\mathfrak U_{k-1}(\frac12a'_1+4w')$.

\emph{Factorization.} Integrating out the configuration of scale $k-1$ against the Haar $\delta$-functions gives, exactly,
\[
\begin{aligned}
&\int\prod_b dU(b)\,\delta\big(U^{(k)},V\big)\,F(U)\\
&\qquad=\int\prod_{b'}dV_{k-1}(b')\ \delta\big(\mathfrak a_k(V_{k-1}),V\big)\int\prod_bdU(b)\,\delta\big(U^{(k-1)},V_{k-1}\big)F(U),
\end{aligned}
\]
for every integrable $F$ that vanishes where $U^{(k)}$ is not defined, since $U^{(k)}=\mathfrak a_k(U^{(k-1)})$.

\emph{The region $\mathscr N'$.} Put $\theta_{k-1}=L^{-3}\theta_0$. For $V_{k-1}\in\mathcal K_0$ the pair $(V_{k-1},\theta_{k-1})$ is admissible at scale $k-1$, in the units of step $k-1$: $V_{k-1}\in\mathfrak U_{k-1}(a_1)$, and by Theorem~\ref{thm:balaban}(b) its minimal orbit lies in $\mathfrak U(B_3(\frac12a'_1+4w'))\subset\mathfrak U(\frac34\theta_{k-1})$, since $B_3a'_1\le L^{-3}\theta_0$. Take $F=\one_{\mathscr N'}e^{-S/g^2}$, with the region $\mathscr N'$ of configurations $U$ for which $V_{k-1}=U^{(k-1)}$ lies in $\mathscr O$ and $U\in\mathscr N_{r'}(V_{k-1})$, the neighbourhood of the setting at scale $k-1$, with $\theta_{k-1}$ in place of $\theta$ and $0<r'\le r_0$. This region has two properties:
\begin{itemize}
\item[(N1)] its intersection with $\{U:U^{(k)}=V\}$ lies in $\mathfrak B(V)\cap\mathfrak U(\theta_0)$;
\item[(N2)] it is a Borel set, and it contains $\mathscr N_r(V)$ for $r$ small.
\end{itemize}
For (N1), let $U\in\mathscr N'$. Then $U\in\mathfrak U(\theta_{k-1})$ in the units of step $k-1$. Passing to the units of step $k$ multiplies curvatures by $L^2$ and currents by $L^3$ (proof of Lemma~\ref{lem:units}). So $U\in\mathfrak U(\theta_0)\subset\mathfrak U_0(c_1)$ in the units of step $k$, and $U\in\mathfrak B(V)$ if $U^{(k)}=V$. This is the reason for the choice of $\theta_{k-1}$.

For (N2), let $\varrho_{k-1}$ be the function of Lemma~\ref{lem:ol-chart-covariance}(c) at scale $k-1$, for $\theta_{k-1}$ and $\mathcal K_0$. It is continuous on the compact set $\mathcal D=\{U\in\mathfrak U(\theta_{k-1}):U^{(k-1)}\in\mathcal K_0\}$, and $\mathscr N'=\{U\in\mathcal D:U^{(k-1)}\in\mathscr O,\ \varrho_{k-1}(U)<r'\}$, by that lemma. This set is relatively open in $\mathcal D$, hence a Borel set. Let $\mathcal D^\circ$ be the set of the $U$ whose curvature and current, in the units of step $k-1$, are smaller than $\theta_{k-1}$, and for which $U^{(k-1)}\in\mathscr O$. It is open, since $\mathscr O$ is open and curvature, current and $U\mapsto U^{(k-1)}$ are continuous there, and $\mathcal D^\circ\subset\mathcal D$. So $\mathscr N^\circ=\{U\in\mathcal D^\circ:\varrho_{k-1}(U)<r'\}$ is open, and $\mathscr N^\circ\subset\mathscr N'$. It contains $\mathscr M(V)$. Indeed, let $u\in\mathfrak G_0$, with restriction $\tilde u\in\mathcal G'$ to the lattice of scale $k-1$. By Lemma~\ref{lem:composition}(b) and Theorem~\ref{thm:variational}(g), $U_k(V)^u$ lies in the minimal orbit, at scale $k-1$, of $(U_k(V)^u)^{(k-1)}=(V^0)^{\tilde u}\in\mathscr O$. So $\varrho_{k-1}(U_k(V)^u)=0$, and $U_k(V)^u\in\mathfrak U(\frac34\theta_{k-1})$ in the units of step $k-1$, by the admissibility shown above; hence $U_k(V)^u\in\mathscr N^\circ$. By Lemma~\ref{lem:ol-chart-covariance}(c) at scale $k$, with $\{V\}$ in place of $\mathcal K_0$, $\mathscr N_r(V)\subset\mathscr N^\circ\subset\mathscr N'$ for $r$ small.

The integrand on the left side of the factorization vanishes unless $U^{(k)}=V$. By (N1), (N2) and Lemma~\ref{lem:ol-neighbourhood}, applied to $\mathscr N'\cap\mathfrak B(V)$, the left side is $I_k(V;g,r)\,(1+O(e^{-c/g^2}))$. On the right side the inner integral is $I_{k-1}(V_{k-1};g,r')$ if $V_{k-1}\in\mathscr O$, and $0$ otherwise. So, up to a factor $1+O(e^{-c/g^2})$, $I_k(V;g,r)$ is the outer integral $T\big(\one_{\mathscr O}\,I_{k-1}(\cdot\,;g,r')\big)(V)$, where $T$ is the map \eqref{eq:rg-map} of the step.

\emph{The inner integral.} The orbit of $V^0$ under $\mathcal G'$ is compact and consists of small-field configurations of scale $k-1$, since $U_{k-1}(V^0)=U_k(V)$ (Lemma~\ref{lem:composition}(b)) and the conditions of the setting are invariant under gauge transformations; this is the first claim. Let $\mathcal K_w$ be the $\mathcal G'$-saturation of the set of the $e^{\ii Z}W$ with $|Z(b)|\le w$ on all bonds. As for $\mathcal K_0$, with $w$ in place of $w'$, it is compact, it contains $\mathcal K_0\supset\mathscr O$ and the closed $w$-neighbourhood of $V^0$, bondwise, and $\mathcal K_w\subset\mathfrak U_{k-1}(\frac12a'_1+4w)$. By the choice of $w$, $8w\le a'_1$ and $16B_3w\le L^{-3}\theta_0$. Then, as in the paragraph ``The region $\mathscr N'$'', $(V_{k-1},\theta_{k-1})$ is admissible at scale $k-1$ for every $V_{k-1}\in\mathcal K_w$, and $V_{k-1}\in\mathfrak U_{k-1}(a_1)$. By Proposition~\ref{prop:ol-laplace} at scale $k-1$, with $\theta_{k-1}$ in place of $\theta$, which does not change $\Gamma^{(1)}_{k-1}$, the expansion of $I_{k-1}$ holds locally uniformly on the set of these configurations, hence uniformly on the compact set $\mathcal K_w$:
\[
I_{k-1}(V_{k-1};g,r')=c_{k-1}(g)\exp\Big(-\frac1{g^2}S\big(U_{k-1}(V_{k-1})\big)+\Gamma^{(1)}_{k-1}(V_{k-1})+o(1)\Big).
\]
This uniformity is the paragraph ``Uniformity in $V$'' of the proof of Proposition~\ref{prop:ol-laplace}, which uses only that $(V_{k-1},\theta_{k-1})$ is admissible and $V_{k-1}\in\mathfrak U_{k-1}(a_1)$: near a point $V^\circ$ of $\mathcal K_w$, with a reference $U_0\in\mathscr M(V^\circ)$, Corollary~\ref{cor:analytic} writes the minimal orbit of $e^{\ii B}V^\circ$ as that of $\mathsf U_{U_0}(A_*(B))$, so the phase, the constraint and the density of Laplace's method depend analytically on $B$, with the Hessian bound uniform. Let $\rho(V_{k-1})=\exp(-g^{-2}S(U_{k-1}(V_{k-1}))+\Gamma^{(1)}_{k-1}(V_{k-1}))$ on $\mathcal K_w$. Since $T$ is positive, $T(\one_{\mathscr O}I_{k-1})=c_{k-1}(g)\,e^{o(1)}\,T(\one_{\mathscr O}\rho)$.

\emph{Gauge fixing.} The function $\rho$ is gauge invariant: $S(U_{k-1}(\cdot))$ by Theorem~\ref{thm:balaban}, $\Gamma^{(1)}_{k-1}$ by Proposition~\ref{prop:ol-laplace}. So $\one_{\mathscr O}\rho$ is $\mathcal G'$-invariant, bounded, and vanishes outside $\mathfrak M$. By Lemma~\ref{lem:gf-step-invariance}(iv), $T(\one_{\mathscr O}\rho)=T^{\rm gf}(\one_{\mathscr O}\rho)$.

\emph{From the region to the cutoff.} We show that
\[
T^{\rm gf}(\one_{\mathscr O}\rho)(V)=T^{\rm gf}(\chi^{\rm w}\rho)(V)\,\big(1+O(\eps_1+\eps_2)\big),\qquad\eps_1=C|\T|e^{-cw^2/g^2},\quad\eps_2=C_\T e^{-c'_Lw^2/g^2},
\]
with $c>0$ absolute, $C$ depending on $G$, $c'_L>0$ depending on $L$, and $C_\T$ depending on $L$ and on the torus; this is the two-radius estimate, with the inner radius $w'$ and the outer radius $w$. Write $\one_{\mathscr O}-\chi^{\rm w}=\one_{\mathscr O}(1-\chi^{\rm w})-(1-\one_{\mathscr O})\chi^{\rm w}$, and let $R_1$, $R_2$ be the two corresponding terms of $T^{\rm gf}(\one_{\mathscr O}\rho)-T^{\rm gf}(\chi^{\rm w}\rho)$. It suffices to show $0\le R_1\le\eps_1T^{\rm gf}(\one_{\mathscr O}\rho)$ and $0\le R_2\le\eps_2T^{\rm gf}(\chi^{\rm w}\rho)$.

\emph{The term $R_1$: gauge translates leaving the cutoff.} Put
\[
h(V_{k-1})=\prod_c\delta\big(\mathfrak a_k(V_{k-1})(c)V(c)^{-1}\big)\,\one_{\mathscr O}(V_{k-1})\,\rho(V_{k-1}).
\]
It is $\mathcal G'$-invariant, by Lemma~\ref{lem:one-step-basic}(b) (with the convention on $\delta$-functions of Section~\ref{sec:fluctuations}). As in the proof of Lemma~\ref{lem:exp-gauge}(a), substituting $V_{k-1}^u$ for $V_{k-1}$, $u\in\mathcal G'$, and averaging over $u$ gives
\[
R_1=\int\prod_bdV_{k-1}(b)\ h\,\psi,\qquad\psi(V_{k-1})=\int_{\mathcal G'}du\ \mathrm{gf}_W(V_{k-1}^u)\big(1-\chi^{\rm w}(V_{k-1}^u)\big),
\]
and, in the same way, $T^{\rm gf}(\one_{\mathscr O}\rho)=\int\prod_bdV_{k-1}(b)\,h$, since $\int_{\mathcal G'}\mathrm{gf}_W(V_{k-1}^u)\,du=1$ on $\mathfrak M$. The function $\psi$ is $\mathcal G'$-invariant, so it suffices to bound it on $\mathscr A$. Let $V_{k-1}\in\mathscr A$, so that $V_{k-1}(b)=e^{\ii Z(b)}W(b)$ with $Z(b)\in\mathfrak g$, $|Z(b)|<w'$, and let $N=\{e^{\ii K}:K\in\mathfrak g,\ |K|\le w/4\}$.
\begin{itemize}
\item If $u(x)\in N$ for all $x$, then $\chi^{\rm w}(V_{k-1}^u)=1$. Indeed,
\[
V_{k-1}^u(b)W(b)^{-1}=u(b_-)\,e^{\ii Z(b)}\,\big(W(b)u(b_+)^{-1}W(b)^{-1}\big),
\]
and by Lemma~\ref{lem:log-estimates}(b),(d) its logarithm is in $\mathfrak g$ and bounded by $\frac w4+w'+\frac w4<w$.
\item Let $x\in B'(y)\setminus\{y\}$, $y\in\T'$, and $u(x)\notin N$. By \eqref{eq:avg-contour-gauge}, $V_{k-1}^u(y,x)W(y,x)^{-1}=e^{\ii H}v^{-1}$ with $e^{\ii H}=V_{k-1}(y,x)W(y,x)^{-1}$ and $v=W(y,x)u(x)W(y,x)^{-1}\notin N$. Since $w'\le c_0L^{-1}$, the proof of Lemma~\ref{lem:exp-gauge}(c) gives $|H|\le4C_MLw'=w/32$. For $v\in G\setminus N$ we have $|v-1|>w/8$: either $|v-1|>c_G\ge w/8$, since $w\le\frac12c_G$ (Definition~\ref{def:chi-w}), or $v=e^{\ii\ell(v)}$ with $\ell(v)\in\mathfrak g$, $|\ell(v)|>w/4$ and $|\ell(v)|\le2|v-1|$ (Lemma~\ref{lem:log-estimates}(a)). Hence $|e^{\ii H}v^{-1}-1|=|v^{-1}-e^{-\ii H}|\ge|v-1|-2|H|\ge w/16$, and the factor of $(y,x)$ in $\mathrm{gf}_W(V_{k-1}^u)$ is at most $z^{-1}e^{-w^2/(512g^2)}$.
\end{itemize}
Since $\chi^{\rm w}$ takes only the values $0$ and $1$ (Definition~\ref{def:chi-w}), $1-\chi^{\rm w}(V_{k-1}^u)\le\sum_x\one(u(x)\notin N)$. The weight $\mathrm{gf}_W(V_{k-1}^u)$ is a product over the points $x$, and each factor integrates to $1$ over $u(x)\in G$. Hence $\psi\le|\T|z^{-1}e^{-w^2/(512g^2)}$ on $\mathscr A$. Since $z(g)^{-1}\le Cg^{-\dim\mathfrak g}$, this is at most $\eps_1$ for $g$ small.

\emph{The term $R_2$: orbits off the inner region.} By Proposition~\ref{prop:exact-step}, applied to the bounded function $(1-\one_{\mathscr O})\rho$,
\[
R_2=z^{-|\T\setminus\T'|}\int\prod_{b\in\mathfrak F}\sigma(Z(b))\,dZ(b)\ \chi^{\rm w}(\hat Z)\prod_c\mathsf J_c\ e^{-g^{-2}F(Z)}\,\big(e^{\Gamma^{(1)}_{k-1}}(1-\one_{\mathscr O})\big)\big(e^{\ii\hat Z(Z)}W\big),
\]
with $F=\mathcal S\circ\hat Z+\Phi_W$ as in Lemma~\ref{lem:step-quadratic}. On the support, $|\hat Z(b)|\le w\le r_2$ on all bonds. By Lemma~\ref{lem:step-quadratic}, Corollary~\ref{cor:covariance} and $C_Lw\le c_L/16$,
\[
F(Z)\ge S(U')+\tfrac18c_L\|Z\|^2-C_Lw\|Z\|^2\ge S(U')+\tfrac1{16}c_L\|Z\|^2 .
\]
Since $e^{\ii\hat Z}W\notin\mathscr A$ and $\hat Z$ is real, some bond has $|\hat Z(b)|\ge w'$; with $|\hat Z|_\infty\le CL^d|Z|_\infty$ (Lemma~\ref{lem:eliminate}(c)) this gives $\|Z\|\ge w'':=w'/(CL^d)$. The factors $\sigma$, $\mathsf J_c$ and $e^{\Gamma^{(1)}_{k-1}}$ are bounded on the domain by a constant depending on $L$ and the torus. With $n=|\mathfrak F|\dim\mathfrak g$ and $\int_{\|Z\|\ge w''}e^{-a\|Z\|^2}dZ\le e^{-aw''^2/2}(2\pi/a)^{n/2}$ for $a=c_L/(16g^2)$, we get
\[
R_2\le C_\T\,z^{-|\T\setminus\T'|}(2\pi g^2)^{n/2}e^{-S(U')/g^2}e^{-c_Lw''^2/(32g^2)} .
\]
By the next step, $T^{\rm gf}(\chi^{\rm w}\rho)\ge\frac12\mathcal N_{k-1}e^{-S(U')/g^2+\Gamma^{\rm step}_{k-1}(U')+\Gamma^{(1)}_{k-1}(V^0)}$ for $g$ small, and $\mathcal N_{k-1}$ is $z^{-|\T\setminus\T'|}(2\pi g^2)^{n/2}$ times a constant (Theorem~\ref{thm:step}(a)). So $R_2\le\eps_2T^{\rm gf}(\chi^{\rm w}\rho)$, with $c'_L=c_L/(32C^2C_1^2L^{2d+2})$.

\emph{The outer integral.} By Lemma~\ref{lem:composition}, the minimal configuration of the step is $V^0$, and $U_{k-1}(V^0)=U_k(V)$. The density $\rho$ has the form \eqref{eq:rho-form} with $\mathfrak E=\Gamma^{(1)}_{k-1}$, real analytic and gauge invariant, on the support of $\chi^{\rm w}$, which is all that Theorem~\ref{thm:step} uses. Theorem~\ref{thm:step}, applied to the step from scale $k-1$ to scale $k$, with $w$ fixed and $g\to0$, gives
\[
\begin{aligned}
T^{\rm gf}(\chi^{\rm w}\rho)(V)={}&\mathcal N_{k-1}\exp\Big(-\frac1{g^2}S\big(U_k(V)\big)+\Gamma^{\rm step}_{k-1}\big(U_k(V)\big)\Big)\\
&\times\int d\mu_{C}(Y)\,\chi^{\rm w}(\hat Z(gY))\,e^{-g^{-2}(\mathcal V+G_3)(gY)+\mathfrak J(gY)+\mathfrak E(e^{\ii\hat Z(gY)}W)} .
\end{aligned}
\]
As $g\to0$:
\begin{itemize}
\item $g^{-2}|(\mathcal V+G_3)(gY)|\le C_Lg|Y|_\infty\|Y\|^2$ (Theorem~\ref{thm:step}(d));
\item $\mathfrak J(gY)=O(g|Y|_\infty)$ (Theorem~\ref{thm:step}(e));
\item $\mathfrak E(e^{\ii\hat Z(gY)}W)\to\mathfrak E(W)=\Gamma^{(1)}_{k-1}(V^0)$, by continuity.
\end{itemize}
On the support of $\chi^{\rm w}$ the integrand is bounded by $e^{C_Lw\|Y\|^2+C}$, which is integrable against $d\mu_C$, since $C\le4c_L^{-1}$ and $C_Lw\le c_L/16$. By dominated convergence the integral tends to $e^{\Gamma^{(1)}_{k-1}(V^0)}$.

\emph{Conclusion.} Collecting the steps,
\[
I_k(V;g,r)=c_{k-1}(g)\mathcal N_{k-1}\exp\Big(-\frac1{g^2}S\big(U_k(V)\big)+\Gamma^{(1)}_{k-1}(V^0)+\Gamma^{\rm step}_{k-1}\big(U_k(V)\big)+o(1)\Big).
\]
Compare with Proposition~\ref{prop:ol-laplace} at scale $k$. The difference $D(V)=\Gamma^{(1)}_k(V)-\Gamma^{(1)}_{k-1}(V^0)-\Gamma^{\rm step}_{k-1}(U_k(V))$ does not depend on $g$, and $D(V)=\log\big(c_{k-1}(g)\mathcal N_{k-1}/c_k(g)\big)+o(1)$ as $g\to0$, where the first term does not depend on $V$. So $D(V)$ is the limit of this term, a constant independent of $V$, and $c_k(g)=c_{k-1}(g)\mathcal N_{k-1}$ up to a constant factor. Of $\Gamma^{(1)}_{k-1}$ the proof has used only Proposition~\ref{prop:ol-laplace}, not its relation to the steps below scale $k-1$.
\end{proof}

\begin{proposition}[Composition of the one-loop terms]\label{prop:ol-composition}
For every $k\ge1$ and every small-field $V$, up to a constant independent of $V$,
\[
\Gamma^{(1)}_k(V)=\sum_{j=0}^{k-1}\Gamma^{\rm step}_j\big(U_k(V)\big).
\]
\end{proposition}

\begin{proof}\hasdetail{prop:ol-composition}
For $k=1$ the proof of Lemma~\ref{lem:ol-one-step} applies to $I_1(V;g,r)=T\big(\one_{\mathscr N_r(V)}e^{-S/g^2}\big)(V)$, with $\rho=e^{-S/g^2}$, i.e.\ with $\mathfrak E=0$ and $\Gamma^{(1)}_0:=0$, and with Lemma~\ref{lem:first-step-formula} in place of Lemma~\ref{lem:composition} and Theorem~\ref{thm:step}. Its part (a) replaces Lemma~\ref{lem:composition}; its parts (b), (c) replace Lemma~\ref{lem:step-quadratic} and Corollary~\ref{cor:covariance}, with $c^{(0)}_L=2C_{\rm P}^{-1}$ in place of $c_L$; and its part (d) replaces Theorem~\ref{thm:step}. Accordingly $r_2$, $C_L$ and $a_{\rm step}$ are replaced by $r^{(0)}_2$, the $C_L$ of that lemma and $a^{(0)}_{\rm step}$. At scale $0$ the fine lattice is the lattice of scale $0$, the base is $W=U_1(V)$ (Lemma~\ref{lem:first-step-formula}(a)), and the region $\mathscr N'$ is replaced by the open set $\mathscr O$ itself. It has the properties (N1) and (N2) of that proof, with $\mathscr O$ in place of $\mathscr N'$.
\begin{itemize}
\item[(N1)] Let $V_0\in\mathscr O$. By Lemma~\ref{lem:first-step-formula}(a) and the choice of its constants, $W\in\mathfrak U_0(L^{-2}(C_*+1)\eps'_1)\subset\mathfrak U_0(\frac14a'_1)$. By the argument before \eqref{eq:chi-w-support}, with $w'$ in place of $w$, $V_0\in\mathfrak U_0(\frac14a'_1+4w')\subset\mathfrak U_0(a'_1)$. As in the proof of Lemma~\ref{lem:first-step-formula}(a), Theorem~\ref{thm:balaban}(b) at $k=0$ gives $V_0\in\mathfrak U(B_3a'_1)$ in the units of step $0$, so its curvature and current are at most $L^3B_3a'_1\le\theta_0$ in the units of step $1$. So $V_0\in\mathfrak B(V)\cap\mathfrak U(\theta_0)$ if $V_0^{(1)}=V$.
\item[(N2)] At $k=1$, $\mathfrak G_0=\mathcal G'$, so $\mathscr M(V)=\{W^u:u\in\mathcal G'\}$, which lies in the open set $\mathscr O$. By Lemma~\ref{lem:ol-chart-covariance}(c), with $\{V\}$ in place of $\mathcal K_0$, $\mathscr N_r(V)\subset\mathscr O$ for $r$ small.
\end{itemize}
By Lemma~\ref{lem:ol-neighbourhood}, applied to $\mathscr O\cap\mathfrak B(V)$, $I_1(V;g,r)=T(\one_{\mathscr O}e^{-S/g^2})(V)\,(1+O(e^{-c/g^2}))$. Proceed as from the step \emph{Gauge fixing} on. This gives $\Gamma^{(1)}_1(V)=\Gamma^{\rm step}_0(U_1(V))$ up to a constant.

Let $k\ge2$ and assume the claim for $k-1$. By Lemma~\ref{lem:ol-one-step}, $V^0=U_k(V)^{(k-1)}$ is small field at scale $k-1$ and $\Gamma^{(1)}_k(V)=\Gamma^{(1)}_{k-1}(V^0)+\Gamma^{\rm step}_{k-1}(U_k(V))$ up to a constant. By the induction hypothesis and Lemma~\ref{lem:composition}(b),
\[
\Gamma^{(1)}_{k-1}(V^0)=\sum_{j<k-1}\Gamma^{\rm step}_j\big(U_{k-1}(V^0)\big)=\sum_{j<k-1}\Gamma^{\rm step}_j\big(U_k(V)\big)
\]
up to a constant. This gives the claim.
\end{proof}

So the one-loop terms of the steps need not be computed one by one. Their sum is the one-loop term of the composite map, a single determinant on the constraint set modulo the residual gauge group. Section~\ref{sec:ol-chart} writes it with the operators of Parts~II and~III.

\subsection{Marginal tensors and lattice symmetry}\label{subsec:ol-marginal}

Let $\Phi$ be a real analytic function of coarse configurations near $1$, on the unit lattice $\T$, invariant under the gauge transformations of $\T$, with vanishing first variation at $1$. Let $\Pi_\Phi(x,y)$ be its second variation at $1$, i.e.\ the kernel of $B\mapsto\frac{d^2}{dt^2}\Phi(e^{\ii tB})|_{t=0}$ with respect to the trace pairing $\sum_b\operatorname{tr}(B(b)B'(b))$ of Section~\ref{sec:action}, with values in $\mathfrak g\otimes\mathfrak g$. We assume that $\Pi_\Phi$ is invariant under the translations of $\T$ and that, as the side of the torus tends to infinity, its kernels converge with uniform exponential decay, in the sense of Lemma~\ref{lem:ol-folding}(e); we use the limit, the infinite-volume second variation, again denoted $\Pi_\Phi$. We use the \emph{bond-midpoint convention}: the value $B(b)$ sits at the midpoint $m_b$ of $b$, and $\hat B_\mu(p)=\sum_be^{-\ii p\cdot m_b}B(b)$, the sum over the positively oriented bonds of direction $\mu$. The Fourier transform $\Pi_\Phi(p)_{\mu\nu}$ is analytic near $p=0$. In the convention with $B(b)$ at $b_-$ it is multiplied by the phases $e^{\ii(p_\mu-p_\nu)/2}$; if $\Pi_\Phi(p)=O(|p|^2)$ in one convention, this changes it by $O(|p|^3)$, so the tensor below does not depend on the convention. In this part $G$ is simple (Section~\ref{sec:intro}, Scope). By $\operatorname{Ad}$-invariance $\Pi_\Phi$ is then proportional to the identity on $\mathfrak g$ \cite[(4.33)]{BalYM4I}, and we suppress the Lie algebra indices relative to the trace pairing.

\begin{definition}[Marginal tensor]\label{def:marginal-tensor}
If $\Pi_\Phi(p)=O(|p|^2)$ as $p\to0$, the \emph{marginal tensor} of $\Phi$ is the tensor $T_\Phi$, symmetric in $\kappa,\lambda$, with $\Pi_\Phi(p)_{\mu\nu}=\sum_{\kappa\lambda}(T_\Phi)_{\mu\nu\kappa\lambda}p_\kappa p_\lambda+O(|p|^3)$. Its \emph{isotropic coefficient} $\beta_\Phi$ is the coefficient of the orthogonal projection of $T_\Phi$, in the Euclidean inner product of tensors, onto the tensor of $\delta_{\mu\nu}|p|^2-p_\mu p_\nu$. With the hypercubic symmetry, $T_\Phi(p)=\beta_\Phi(\delta_{\mu\nu}|p|^2-p_\mu p_\nu)$ (Lemma~\ref{lem:ol-symmetry}(c)), and $\beta_\Phi$ is Ba\l aban's coefficient \cite[(5.37), (5.42)]{BalYM4I}.
\end{definition}

\begin{remark}[The orders $0$ and $1$ in $p$]\label{rem:marginal-orders}
Since $\Phi$ is gauge invariant and its first variation at $1$ vanishes, $\Pi_\Phi(B,D\lambda)=0$ for all $B$ and scalar fields $\lambda$. In the bond-midpoint convention $\widehat{D\lambda}_\mu(p)=\ii\hat p_\mu\hat\lambda(p)$ with $\hat p_\mu=2\sin(p_\mu/2)$, so $\Pi_\Phi$ is transverse: $\sum_\mu\hat p_\mu\Pi_\Phi(p)_{\mu\nu}=0$. In infinite volume a constant field $B$ is pure gauge to first order, so $\Pi_\Phi(0)=0$. A term linear in $p$ is not excluded by transversality and reality: $\ii c\,\varepsilon_{\mu\nu\lambda\sigma}n_\sigma p_\lambda$, with a lattice vector $n$, is of Chern--Simons type. It is excluded by the point reflection $x\mapsto-x$, which preserves orientation in $d=4$. The functions of this part are invariant under it, by Lemma~\ref{lem:ol-symmetry} below. A nonzero linear term would not be harmless: under the change of scale it would produce relevant terms of order $1$ per step.
\end{remark}

Let $\mathfrak S_k$ be the group of lattice symmetries of Section~\ref{sec:setup}, generated by the translations of $\T$ and by $B_4$ about the points of $\T$. It acts on configurations by $(rV)(b)=V(r^{-1}b)$.

\begin{lemma}[Lattice symmetry of the one-loop terms]\label{lem:ol-symmetry}
\begin{enumerate}[label=(\alph*)]
\item Let $V$ be small field and $r\in\mathfrak S_k$. Then $rV$ is small field, and $\Gamma^{(1)}_k(rV)=\Gamma^{(1)}_k(V)$.
\item Let $j\ge0$, and let $V'$ be a small-field configuration of scale $j+1$, in the units of step $j+1$. Then $\Gamma^{\rm step}_j(U_{j+1}(rV'))=\Gamma^{\rm step}_j(U_{j+1}(V'))$ for every $r\in\mathfrak S_{j+1}$. Here, as everywhere in this section, $\Gamma^{\rm step}_j$ is evaluated at the background fields, the critical references of Theorem~\ref{thm:step}.
\item Let $\Phi$ be one of the functions $\Gamma^{(1)}_k$ or $V'\mapsto\Gamma^{\rm step}_j(U_{j+1}(V'))$ with $j\ge0$. Assume that its second variations at $1$ on the tori converge with uniform exponential decay as the side tends to infinity (Lemma~\ref{lem:ol-folding}(e)). Then, in the bond-midpoint convention, $\Pi_\Phi(Rp)=R\,\Pi_\Phi(p)R^T$ for $R\in B_4$; in particular $\Pi_\Phi(-p)=\Pi_\Phi(p)$. Hence $\Pi_\Phi(p)=O(|p|^2)$, the hypothesis of Definition~\ref{def:marginal-tensor} holds, and
\[
T_\Phi(p)_{\mu\nu}=\beta_\Phi\big(\delta_{\mu\nu}|p|^2-p_\mu p_\nu\big)
\]
with a single coefficient $\beta_\Phi$.
\end{enumerate}
\end{lemma}

The hypothesis of (c) is proved in Lemma~\ref{lem:ol-torus-limit}(c) for $\Gamma^{(1)}_k$, and in Lemma~\ref{lem:ol-limits} for the functions $\Gamma^{\rm step}_j$.

\begin{proof}\hasdetail{lem:ol-symmetry}
(a) By Theorem~\ref{thm:balaban}(c), $rV$ has the same curvature as $V$, and its minimal orbit is the image under $r$ of that of $V$. By Lemma~\ref{lem:one-step-basic}, $U_k(rV)^{(j+1)}$ is a gauge transform of $r(U_k(V)^{(j+1)})$ for every $j<k$. The conditions of the setting are bounds on curvatures and on the side of the torus, which are invariant under gauge transformations and lattice symmetries. So $rV$ is small field. The class $\mathfrak U(\theta')$ is invariant under $\mathfrak S_k$ (proof of Theorem~\ref{thm:variational}(h)), so the same $\theta$ serves for $V$ and for $rV$.

To avoid a clash with the symmetry $r$, write $r'\in(0,r_0]$ for the radius of the neighbourhoods, and let $\mathscr N=\mathscr N_{r'}(V)\cup r^{-1}\mathscr N_{r'}(rV)$. The map $U\mapsto rU$ maps $\mathfrak B(V)\cap\mathfrak U(\theta_0)$ onto $\mathfrak B(rV)\cap\mathfrak U(\theta_0)$ and $\mathfrak G_0$-orbits onto $\mathfrak G_0$-orbits (proof of Theorem~\ref{thm:balaban}(c)), and it is a homeomorphism. So $\mathscr N$ is a Borel subset of $\mathfrak B(V)\cap\mathfrak U(\theta_0)$ that contains $\mathscr N_{r'}(V)$, by Lemma~\ref{lem:ol-chart-covariance}(c), and $r\mathscr N$ is one of $\mathfrak B(rV)\cap\mathfrak U(\theta_0)$ that contains $\mathscr N_{r'}(rV)$. By Lemma~\ref{lem:ol-neighbourhood}, the integrals over $\mathscr N$ at $V$ and over $r\mathscr N$ at $rV$ equal $I_k(V;g,r')$ and $I_k(rV;g,r')$, respectively, up to factors $1+O(e^{-c/g^2})$. Substitute $U=r\tilde U$ in the integral over $r\mathscr N$ at $rV$:
\begin{itemize}
\item the map $\tilde U\mapsto r\tilde U$ permutes the bond variables and inverts some of them, so it preserves $\prod_bdU(b)$;
\item $(r\tilde U)^{(k)}=r\,\tilde U^{(k)}$ by Lemma~\ref{lem:one-step-basic}, and $\delta(r\tilde U^{(k)},rV)=\delta(\tilde U^{(k)},V)$, since $\delta$ is invariant under conjugation and inversion (proof of Proposition~\ref{prop:symmetry}(a));
\item $S(r\tilde U)=S(\tilde U)$, and $\one_{r\mathscr N}(r\tilde U)=\one_{\mathscr N}(\tilde U)$.
\end{itemize}
So the two integrals are equal, and $I_k(rV;g,r')=I_k(V;g,r')(1+O(e^{-c/g^2}))$. By Theorem~\ref{thm:balaban}(c), $S(U_k(rV))=S(U_k(V))$. Proposition~\ref{prop:ol-laplace}, at $V$ and at $rV$, gives $\Gamma^{(1)}_k(rV)-\Gamma^{(1)}_k(V)=o(1)$ as $g\to0$. Both terms are independent of $g$, so they are equal.

(b) We use the notation of Sections~\ref{sec:fluctuations}--\ref{sec:covariance} for the step from scale $j$ to scale $j+1$, in the units of step $j$, with $j$ in place of $k$: $U'=U_{j+1}(V')$, its base $W=U'^{(j)}$, the operator $\Delta_j=\Delta_j(W,U')$ of Corollary~\ref{cor:covariance}, and the Jacobians $\mathfrak j_c=\mathfrak j_c(W)$ of Theorem~\ref{thm:step}(b). For $j=0$, $W=U'$, and $\Delta_0$ and its properties are those of Lemma~\ref{lem:first-step-formula}(a)--(c). So
\[
\Gamma^{\rm step}_j(U')=-\tfrac12\log\det\Delta_j(W,U')+\sum_c\log\mathfrak j_c(W).
\]
This does not depend on the choice of $U'$ in its minimal orbit. Indeed, another choice is $U'^u$, with $u$ real and $u|_{\T'}=1$, and with the base $W^{u|_\T}$. By Corollary~\ref{cor:covariance}(a) (Lemma~\ref{lem:first-step-formula}(c) for $j=0$), $\Delta_j$ is then replaced by $R(u)\Delta_jR(u)^{-1}$, with $R(u)$ orthogonal, and the $\mathfrak j_c$ are invariant (Proposition~\ref{prop:jacobians}(b)).

Let $r\in\mathfrak S_{j+1}$. By (a), $rV'$ is small field. By Proposition~\ref{prop:symmetry}(b), with $j+1$ in place of $k$, $rU'$ is an element of the minimal orbit of $rV'$; its base is $(rU')^{(j)}=rW$, by Lemma~\ref{lem:one-step-basic}, since $\mathfrak S_{j+1}\subset\mathfrak S_j$. By Corollary~\ref{cor:covariance}(b) (Lemma~\ref{lem:first-step-formula}(c) for $j=0$),
\[
\Delta_j(rW,rU')=r_W\,\Delta_j(W,U')\,r_W^{-1},
\]
where $r_W$ maps the real fields on $\mathfrak F$ orthogonally onto themselves. So $\log\det\Delta_j(rW,rU')=\log\det\Delta_j(W,U')$; equivalently, this is Proposition~\ref{prop:det-densities}(d), summed over $b\in\mathfrak F$, since $r$ maps $\mathfrak F$ onto $\mathfrak F$. By Proposition~\ref{prop:jacobians}(d), $\sum_c\log\mathfrak j_c(rW)=\sum_c\log\mathfrak j_c(W)$. Hence $\Gamma^{\rm step}_j(U_{j+1}(rV'))=\Gamma^{\rm step}_j(rU')=\Gamma^{\rm step}_j(U')$.

(c) The configuration $V=1$ is small field. The function $\Phi$ is real analytic and gauge invariant. For $\Gamma^{(1)}_k$ this is Proposition~\ref{prop:ol-laplace}. For $V'\mapsto\Gamma^{\rm step}_j(U_{j+1}(V'))$ it follows from the identity $\Gamma^{\rm step}_j(U_{j+1}(V'))=\Gamma^{(1)}_{j+1}(V')-\Gamma^{(1)}_j(V_j(V'))+\kappa_j$ of the paragraph after this proof, with $V_j(V')=U_{j+1}(V')^{(j)}$ and $\Gamma^{(1)}_0=0$; this identity uses Lemma~\ref{lem:ol-one-step} and Proposition~\ref{prop:ol-composition}, not this lemma. Indeed, $V_j$ is real analytic up to gauge transformations of $\T^{(j)}$ (Corollary~\ref{cor:analytic}), on which $\Gamma^{(1)}_j$ does not depend, and $V_j(V'^{v})$ is a gauge transform of $V_j(V')$. Its first variation at $1$ vanishes, since $G$ is semisimple (Section~\ref{sec:rg-step}); so it satisfies the assumptions made above Definition~\ref{def:marginal-tensor}, with the translation invariance given by (a) and (b). For $r(x)=Rx+v$ in $\mathfrak S_k$ (respectively $\mathfrak S_{j+1}$) and real $B$, $r(e^{\ii tB})=e^{\ii t\,r_1B}$, where $(r_1B)(b)=B(r^{-1}b)$ and $B(-b)=-B(b)$. By (a) and (b), $\Pi_\Phi(r_1B,r_1B)=\Pi_\Phi(B,B)$. The map $r$ sends the bond with midpoint $m$ and direction $e_\nu$ to the bond with midpoint $r(m)$ and direction $Re_\nu=\pm e_{\sigma(\nu)}$, reversed if the sign is negative. Hence, regarding $B$ as a vector field on the midpoints, $(r_1B)(r(m))=R\,B(m)$, and $\widehat{r_1B}(p)=e^{-\ii p\cdot v}R\,\hat B(R^{-1}p)$. Since $R$ is orthogonal, the invariance of the form, extended sesquilinearly to complex fields, is $\Pi_\Phi(Rp)=R\,\Pi_\Phi(p)R^T$. For $R=-1$ this is $\Pi_\Phi(-p)=\Pi_\Phi(p)$. The kernel of $\Pi_\Phi$ is real and symmetric, so $\Pi_\Phi(p)^T=\Pi_\Phi(-p)=\overline{\Pi_\Phi(p)}$; with the evenness, $\Pi_\Phi(p)$ is real and symmetric.

Expand $\Pi_\Phi(p)=\Pi^{(1)}(p)+\Pi^{(2)}(p)+O(|p|^3)$ with $\Pi^{(i)}$ homogeneous of degree $i$; there is no constant term, since $\Pi_\Phi(0)=0$ (Remark~\ref{rem:marginal-orders}). Evenness gives $\Pi^{(1)}=0$, so $\Pi_\Phi(p)=O(|p|^2)$ and $T_\Phi=\Pi^{(2)}$. Write $T_\Phi(p)_{\mu\nu}=\sum_{\kappa\lambda}T_{\mu\nu\kappa\lambda}p_\kappa p_\lambda$ with $T_{\mu\nu\kappa\lambda}$ real and symmetric in $\kappa,\lambda$. Invariance under the reflections of single axes shows that $T_{\mu\nu\kappa\lambda}=0$ unless each value occurs an even number of times among $\mu,\nu,\kappa,\lambda$; invariance under the permutations of the axes shows that the remaining coefficients depend only on the pattern of equal indices. Hence
\[
T_\Phi(p)_{\mu\nu}=t_1\,\delta_{\mu\nu}|p|^2+t_2\,p_\mu p_\nu+t_3\,\delta_{\mu\nu}p_\mu^2
\]
for real $t_1,t_2,t_3$. The part of order $|p|^3$ of the transversality $\sum_\mu\hat p_\mu\Pi_\Phi(p)_{\mu\nu}=0$ is $\sum_\mu p_\mu T_\Phi(p)_{\mu\nu}=(t_1+t_2)|p|^2p_\nu+t_3\,p_\nu^3=0$ for all $p$. So $t_3=0$ and $t_2=-t_1$, and $\beta_\Phi=t_1$.
\end{proof}

The proof of (a) uses only the invariance of the Haar measure and of the $\delta$-function, $S\circ r=S$, the covariance of the averaging (Lemma~\ref{lem:one-step-basic}(d)), the covariance of the minimal orbits (Theorem~\ref{thm:balaban}(c), Proposition~\ref{prop:symmetry}(b)) and the freedom in the neighbourhood (Lemma~\ref{lem:ol-neighbourhood}); it uses no gauge fixing. The proof of (b) uses, besides the covariance of the minimal orbits, the lattice covariance of the coarse action (Proposition~\ref{prop:hessian}(e)), of the elimination, and of the quadratic part $\|\mathfrak G_W\cdot\|^2$ of the gauge fixing (Corollary~\ref{cor:covariance}(b)); the last is that of the averaged contour variables (Lemma~\ref{lem:avg-transport}(d)). Part~(b) also follows from (a), by a route that uses no covariance property of the gauge fixing. For $j\ge1$, Lemma~\ref{lem:ol-one-step} at $k=j+1$ gives $\Gamma^{\rm step}_j(U_{j+1}(V'))=\Gamma^{(1)}_{j+1}(V')-\Gamma^{(1)}_j(V_j(V'))+\kappa_j$, with $V_j(V')=U_{j+1}(V')^{(j)}$ and a constant $\kappa_j$; configurations are carried from the units of one step to those of another without change (Section~\ref{sec:rg-step}). By Theorem~\ref{thm:balaban}(c) and Lemma~\ref{lem:one-step-basic}, $V_j(rV')$ is a gauge transform of $rV_j(V')$, and (a) at the scales $j+1$ and $j$, with the gauge invariance of $\Gamma^{(1)}_j$, gives the claim. For $j=0$, Proposition~\ref{prop:ol-composition} at $k=1$ gives $\Gamma^{\rm step}_0(U_1(V'))=\Gamma^{(1)}_1(V')+\kappa_0$. On this route the gauge fixing enters only through the exact identity of Lemma~\ref{lem:gf-step-invariance}(iv) on invariant integrands, through the explicit product form of $\mathrm{gf}_W$ and Lemma~\ref{lem:exp-gauge} in the bound on $R_1$, and through the gauge-fixed formula of Theorem~\ref{thm:step}, and of Lemma~\ref{lem:first-step-formula} for $j=0$, which defines $\Gamma^{\rm step}_j$. Neither route uses the covariance of the gauge-field operators of Parts~II and~III, which are covariant under reflections only at flat backgrounds (Remark~\ref{rem:symmetry-reflections}).

Part~(b) is a statement at the background fields $U'=U_{j+1}(V')$. There $\Delta_j$ is the Hessian of the coarse action, composed with the elimination, plus the gauge-fixing term (Lemma~\ref{lem:step-quadratic}, Proposition~\ref{prop:hessian-general}(a)), and its covariance under the reflections in $\mathfrak S_{j+1}$ is that of the coarse action, Proposition~\ref{prop:hessian}(e), which rests on Theorem~\ref{thm:variational}(h). At references that are not critical, $\Delta_j$ is defined in Section~\ref{subsec:sq-hessian} through $\boldsymbol\Gamma_\pi[W_2]$ and $\bH_0$, which are not shown to be covariant under reflections at curved backgrounds (Remark~\ref{rem:symmetry-reflections}). So no covariance under reflections is claimed for $\Delta_j$, for the pieces $\mathfrak e(X,U')$ of Proposition~\ref{prop:one-loop-polymer}, or for the one-loop part $E'_{\rm ol}$ of the raw new generation, at such references. Nothing below needs it: Section~\ref{sec:ol-uv} uses (b) only through the second variations at $V'=1$, i.e.\ at the background fields $U_{j+1}(V')$ with $V'$ near $1$ (proofs of Theorem~\ref{thm:ol-cumulative} and Proposition~\ref{prop:flow-F}).

The lattice symmetry rests on the symmetry of Ba\l aban's averaging \cite[(0.3)--(0.12)]{BalYM4I}: with blocks that have the coarse point at a corner, or with a single ordered contour per point, Lemma~\ref{lem:one-step-basic}(d) fails. The invariance under single reflections, which reverse the orientation, also excludes a topological term: the density $\varepsilon_{\mu\nu\kappa\lambda}\operatorname{tr}(F_{\mu\nu}F_{\kappa\lambda})$ is odd under them, so it cannot occur in a local expansion of $\Gamma^{(1)}_k$ or of the step functions. (At the level of second variations it is invisible anyway, since its quadratic part is a total derivative.)

\subsection{Periodic operators at \texorpdfstring{$U=1$}{U=1} and infinite volume}\label{subsec:ol-infinite-volume}

Marginal tensors are read off in infinite volume, and the change of scale (Lemma~\ref{lem:ol-pullback}) applies the minimizer map of Part~I to linear fields, which are not periodic. At $U=1$ the operators of Part~I commute with the translations of the unit lattice. This subsection constructs their infinite-volume versions from the operators on the tori, by folding (Lemma~\ref{lem:ol-folding}, Corollary~\ref{cor:ol-part-one}). It transfers their bounds through the Bloch fibres (Lemma~\ref{lem:ol-bloch}), collects the trace per cell used in Section~\ref{sec:ol-vacpol} (Lemma~\ref{lem:ol-trace-tools}), and determines the infinite-volume minimizer on linear fields (Lemma~\ref{lem:ol-liouville}, Proposition~\ref{prop:ol-IM}, Corollary~\ref{cor:ol-umklapp}). Only Part~I is used.

Throughout, $k\ge1$ is fixed, and we use the units of scale $k$: the fine lattice is $\eta\Z^4$, with $\eta=L^{-k}=N^{-1}$, and the unit lattice is $\Z^4\subset\eta\Z^4$. Constants may depend on $\eta$; what matters is that they do not depend on the side $n$ of the torus.

\paragraph{Fields and kernels.} A \emph{field space} consists of the fields on the points, on the positively oriented bonds, or on the positively oriented plaquettes of a lattice $\delta\Z^4$, where $\delta^{-1}$ is a divisor of $N$, with values in a finite-dimensional complex vector space $E$ with a Hermitian product. Examples are $\delta=\eta$ or $\delta=1$, and $E=\C$ or $E=\mathfrak g^c$. These points, bonds or plaquettes are the \emph{sites}, and $\Sigma$ denotes the set of sites. A site $s$ has a base point $x(s)$, the point itself or the initial point of the bond or plaquette, and $|s-s'|=|x(s)-x(s')|$. The translations by $a\in\Z^4$ act on $\Sigma$, $s\mapsto s+a$. Since $\delta^{-1}$ is odd, no base point has a half-integer coordinate, so $\Sigma$ is the disjoint union of the translates of the \emph{cell} $\Sigma_0$, the set of the sites with base point in $\{x:|x_\mu|<\frac12\ \text{for all }\mu\}$. With the weight $w=\delta^d$, the inner product is $\ip fg=\sum_sw\,\ip{f(s)}{g(s)}_E$; on $\eta\Z^4$ and $\Z^4$ it is that of Section~\ref{sec:setup}. A field on the torus of side $n$ is an $n\Z^4$-periodic field. Its inner product is the same sum over the sites of the torus, the classes in $\Sigma/n\Z^4$, and $|s-s'|_n$ is the distance on the torus. The operators of Part~I act componentwise in a basis of $E$. Their kernels are real, so their bounds, proved for real fields, hold for $E$-valued fields with the same constants. A field $f$ has \emph{polynomial growth} if $|f(s)|\le C(1+|x(s)|)^m$ for some $C$ and $m$.

An operator $K$ from the fields on $\Sigma$ to those on $\Sigma'$ has a kernel $K(s,s')\in\operatorname{Hom}(E,E')$, with $(Kf)(s)=\sum_{s'}K(s,s')f(s')$. On a torus the sum runs over the sites of the torus, and for $s,s'\in\Sigma$, $K(s,s')$ denotes the kernel at their classes. The weights enter the inner products, not the kernels. The kernel has \emph{kernel decay} $(C,\gamma)$ if $|K(s,s')|\le Ce^{-\gamma|s-s'|}$, and it is \emph{periodic} if $K(s+a,s'+a)=K(s,s')$ for $a\in\Z^4$. Let $\mathscr P$ be the set of the periodic kernels with kernel decay $(C,\gamma)$ for some $C,\gamma>0$. For $K\in\mathscr P$ and $n\ge1$, the \emph{periodization} $[K]_n$ is the operator on the torus of side $n$ with the kernel
\[
[K]_n(s,s')=\sum_{a\in n\Z^4}K(s,s'+a),
\]
which does not depend on the representatives $s,s'$. Put $\sigma_\gamma=\sup_{t\in\R^4}\sum_se^{-\gamma|x(s)-t|/2}$, the sum over the sites of a field space; it depends on $\gamma$, $\eta$ and the field space.

\paragraph{Local families.} For each $n\ge8$ let $K_n$ be an operator between field spaces on the torus of side $n$. The family $(K_n)$ is \emph{local} if there are a periodic kernel $K$ and an $r$ such that $K(s,s')=0$ for $|s-s'|>r$ and $K_n=[K]_n$ for all $n\ge8$. The operators of Section~\ref{sec:setup}, namely $\partial_\mu$, $\partial^*_\mu$, $d$, $d^*$, $\dl$, $\dl^*$, the directional convolutions, $\iota$, $\iota^*$, $Q_\sigma$, $Q_\sigma^*$, $\bQ$ and $\bQ^*$, are local, and so is the identity. On every torus each of them is given by one formula, with the points taken modulo the torus, and this is the periodization of the same formula on $\eta\Z^4$. A family is \emph{algebraic in local operators} if it is obtained from local families by finitely many sums, multiplications by scalars, products, adjoints and inverses; the inverse $(K_n^{-1})$ is formed only if every $K_n$ is invertible.

\begin{lemma}[Periodic kernels and folding]\label{lem:ol-folding}
\begin{enumerate}[label=(\alph*)]
\item \emph{(Calculus.)} A kernel with kernel decay maps the fields of polynomial growth to fields of polynomial growth, and products of such kernels act on them associatively. If $K_1$ and $K_2$ have kernel decay $(C_1,\gamma)$ and $(C_2,\gamma)$, then $K_1K_2$ has kernel decay $(C_1C_2\sigma_\gamma,\gamma/2)$, with $\sigma_\gamma$ of the intermediate field space. The set $\mathscr P$ is closed under sums, products and adjoints. For $K\in\mathscr P$ and $n\ge1$, $[K]_n$ acts on the $n$-periodic fields as $K$ does; $K\mapsto[K]_n$ respects sums, products and adjoints, and $[I]_n=I$. If $K$ has kernel decay $(C,\gamma)$, then, with $S_\gamma$ of Section~\ref{sec:decay},
\[
\big|[K]_n(s,s')-K(s,s')\big|\le CS_{\gamma/4}\,e^{-\gamma n/2}\qquad\text{for }|s-s'|\le n/4 .
\]
In particular $K=0$ if $[K]_n=0$ for infinitely many $n$.
\item \emph{(Folding.)} Let $(K_n)_{n\ge8}$ be algebraic in local operators. Then $K_n$ commutes with the translations by the vectors of $\Z^4$. For $8\le n\mid n'$, $K_{n'}$ maps the $n$-periodic fields to themselves and acts on them as $K_n$. Equivalently, $K_n(s,s')=\sum_aK_{n'}(s,s'+a)$, the sum over a set of representatives $a$ of $n\Z^4/n'\Z^4$.
\item \emph{(Limit.)} Let $(K_n)_{n\ge8}$ be algebraic in local operators, with kernel decay uniform in $n$: $|K_n(s,s')|\le Ce^{-\gamma|s-s'|_n}$ for all $n\ge8$. Then there is a unique $K_\infty\in\mathscr P$ with $[K_\infty]_n=K_n$ for all $n\ge8$. It has kernel decay $(C,\gamma)$, and $|K_n(s,s')-K_\infty(s,s')|\le CS_{\gamma/4}e^{-\gamma n/2}$ for $|s-s'|\le n/4$.
\item \emph{(Algebra.)} If $(K_n)$ and $(K'_n)$ are as in (c), so are their sums, products and adjoints, and their limits are $K_\infty+K'_\infty$, $K_\infty K'_\infty$ and $K_\infty^*$. If $(K_n)$ and $(K_n^{-1})$ are as in (c), then $(K^{-1})_\infty$ is the inverse of $K_\infty$: $K_\infty(K^{-1})_\infty=(K^{-1})_\infty K_\infty=I$. So every identity between such families that holds on all tori holds between their limits, in $\mathscr P$ and, by (a), on the fields of polynomial growth.
\item \emph{(Limits of products.)} Let $\mathfrak N$ be an infinite set of sides, and for $n\in\mathfrak N$ let $X_n$ be an operator between field spaces on the torus of side $n$. The family $(X_n)$ \emph{converges with uniform decay} $(C,\gamma)$ to a kernel $X_\infty$ on the infinite lattice if $|X_n(s,s')|\le Ce^{-\gamma|s-s'|_n}$ for all $n\in\mathfrak N$, and $X_n(s,s')\to X_\infty(s,s')$ for all $s,s'$ as $n\to\infty$ in $\mathfrak N$. Then $X_\infty$ has kernel decay $(C,\gamma)$. If $(X_n)$ and $(Y_n)$ converge with uniform decay $(C,\gamma)$ to $X_\infty$ and $Y_\infty$, then $(X_nY_n)$ converges with uniform decay $(C^2\sigma_\gamma,\gamma/2)$ to $X_\infty Y_\infty$, and the transposed kernels $X_n^{\sf T}(s',s)=X_n(s,s')^{\sf T}$ converge with uniform decay $(C,\gamma)$ to $X_\infty^{\sf T}$. The families of (c) converge with uniform decay $(C,\gamma)$ to their limits.
\end{enumerate}
\end{lemma}

\begin{proof}\hasdetail{lem:ol-folding}
(a) Let $|K(s,s')|\le Ce^{-\gamma|s-s'|}$ and $|f(s')|\le C'(1+|x(s')|)^m$. Since $1+|x(s')|\le(1+|x(s)|)(1+|s-s'|)$,
\[
\sum_{s'}|K(s,s')f(s')|\le CC'(1+|x(s)|)^m\sum_{s'}(1+|s-s'|)^me^{-\gamma|s-s'|},
\]
and the last sum is bounded uniformly in $s$, since a unit cube contains the base points of at most $2^d|\Sigma_0|$ sites. This bound, applied twice, justifies the exchange of the sums in $(K_1K_2)f=K_1(K_2f)$. For the product, $|s-s'|+|s'-s''|\ge|s-s''|$ gives $\sum_{s'}e^{-\gamma|s-s'|-\gamma|s'-s''|}\le e^{-\gamma|s-s''|/2}\sum_{s'}e^{-\gamma|s-s'|/2}\le\sigma_\gamma e^{-\gamma|s-s''|/2}$. The adjoint of $K$ has the kernel $K^*(s',s)=w(s)w(s')^{-1}K(s,s')^*$. Sums, products and adjoints of periodic kernels are periodic. If $K$ is periodic and $f$ is $n$-periodic, grouping the sites $s'$ by their classes modulo $n\Z^4$ gives $(Kf)(s)=\sum_{s'}[K]_n(s,s')f(s')$, the sum over the sites of the torus. So $[K]_n$ is the restriction of $K$ to the $n$-periodic fields, which gives the statements on sums and products and $[I]_n=I$. For adjoints, the same grouping in $\sum_{s,s'}w(s)\ip{g(s)}{K(s,s')f(s')}_E$, with $s$ over the sites of the torus and $s'\in\Sigma$, gives $\ip g{[K]_nf}=\ip{[K^*]_ng}f$ for $n$-periodic $f$ and $g$. Finally, for $|s-s'|\le n/4$ and $0\ne a\in n\Z^4$ we have $|a|\ge n$ and $|x(s)-x(s')-a|\ge|a|-n/4\ge\frac34|a|$, so
\[
\big|[K]_n(s,s')-K(s,s')\big|\le C\sum_{0\ne a\in n\Z^4}e^{-3\gamma|a|/4}\le Ce^{-\gamma n/2}\sum_{a\in\Z^4}e^{-\gamma|a|/4}=CS_{\gamma/4}e^{-\gamma n/2} .
\]
If $[K]_n=0$ for infinitely many $n$, letting $n\to\infty$ gives $K=0$.

(b) We argue along the construction of the family. A local family $([K]_n)$ commutes with the translations, since $K$ is periodic, and by (a) both $[K]_{n'}$ and $[K]_n$ act on the $n$-periodic fields as $K$. Sums, multiples and products preserve the two properties. The translations are unitary, so $K_{n'}^*$ commutes with them, and it maps the $n$-periodic fields, the common fixed points of the translations by the $ne_\mu$, to themselves. On $n$-periodic fields, the inner product of the torus of side $n'$ is $(n'/n)^d$ times that of the torus of side $n$. So, for $n$-periodic $f$ and $g$, $\ip g{K_{n'}^*f}=(n'/n)^d\ip{K_ng}f=\ip g{K_n^*f}$, both in the inner product of the torus of side $n'$, and hence $K_{n'}^*f=K_n^*f$. If every $K_n$ is invertible and $f$ is $n$-periodic, then $u=K_n^{-1}f$ is $n$-periodic and $K_{n'}u=K_nu=f$. So $K_{n'}^{-1}f=K_n^{-1}f$, and $K_{n'}^{-1}$ commutes with the translations, as $K_{n'}$ does. The identity for the kernels follows by applying $K_{n'}$ to the $n$-periodic field $\sum_a\delta_{s'+a}e$, $e\in E$, which is the field $\delta_{s'}e$ of the torus of side $n$.

(c) Let $s,s'\in\Sigma$ and $8\le n\mid n'$. Let $A(n')$ be the set of the $a\in n\Z^4$ with $x(s)-x(s')-a\in(-\frac{n'}2,\frac{n'}2]^d$. It is a set of representatives of $n\Z^4/n'\Z^4$, and for $a\in A(n')$ the distance on the torus of side $n'$ is $|s-(s'+a)|_{n'}=|x(s)-x(s')-a|$. By (b),
\[
K_n(s,s')=\sum_{a\in A(n')}K_{n'}(s,s'+a).
\]
If $|s-s'|\le n/4$, then $0\in A(n')$, and the terms with $a\ne0$ sum to at most $CS_{\gamma/4}e^{-\gamma n/2}$, as in (a). For arbitrary $n_1,n_2\ge8$ with $|s-s'|\le\min(n_1,n_2)/4$, comparing both with $n'=n_1n_2$ gives $|K_{n_1}(s,s')-K_{n_2}(s,s')|\le2CS_{\gamma/4}e^{-\gamma\min(n_1,n_2)/2}$. So $K_n(s,s')$ converges as $n\to\infty$; let $K_\infty(s,s')$ be the limit. By (b), $K_n(s+a,s'+a)=K_n(s,s')$ for $a\in\Z^4$, so $K_\infty$ is periodic. Since $|s-s'|_n=|s-s'|$ for $n$ large, $|K_\infty(s,s')|\le Ce^{-\gamma|s-s'|}$. Now fix $n\ge8$ and $s,s'$, and let $n'=nm$ with $m\to\infty$. The sets $A(n')$ increase to $n\Z^4$, the terms $K_{n'}(s,s'+a)$ converge to $K_\infty(s,s'+a)$, and they are bounded by $Ce^{-\gamma|x(s)-x(s')-a|}$, which is summable over $a\in n\Z^4$. By dominated convergence, $K_n(s,s')=[K_\infty]_n(s,s')$. The uniqueness and the bound follow from (a).

(d) The sums, products and adjoints are algebraic in local operators. Their kernel decay is uniform in $n$: for products, represent the sites $s'$ of the torus by the $s'\in\Sigma$ with $x(s')-x(s)\in(-\frac n2,\frac n2]^d$, so that $|s-s'|_n=|s-s'|$, and argue as in (a); for adjoints, use the formula for $K^*$ in (a). Since $[\cdot]_n$ respects the operations, $[K_\infty K'_\infty]_n=K_nK'_n$ for all $n\ge8$, so $K_\infty K'_\infty$ is the limit of $(K_nK'_n)$, by the uniqueness in (c); likewise for sums and adjoints. If $(K_n^{-1})$ is as in (c), then $[K_\infty(K^{-1})_\infty]_n=K_nK_n^{-1}=[I]_n$ for all $n$, so $K_\infty(K^{-1})_\infty=I$ by (a), and likewise in the other order.

(e) The bound on $X_\infty$ holds since $|s-s'|_n=|s-s'|$ for $n$ large. For the product, fix $s$ and $s''$, and write $(X_nY_n)(s,s'')=\sum_{s'\in D_n}X_n(s,s')Y_n(s',s'')$, where $D_n$ is the set of the $s'\in\Sigma$ with $x(s')-x(s)\in(-\frac n2,\frac n2]^d$. For $s'\in D_n$ the term is bounded by $C^2e^{-\gamma|s-s'|}$, which is summable over $s'\in\Sigma$; every $s'$ lies in $D_n$ for $n$ large, and the terms converge to $X_\infty(s,s')Y_\infty(s',s'')$. Dominated convergence gives the limit $(X_\infty Y_\infty)(s,s'')$, and the uniform decay follows as in (d). The statement on transposes is clear, and the last one is (c).
\end{proof}

\begin{corollary}[The operators of Part~I in infinite volume]\label{cor:ol-part-one}
The following families, on the tori of all sides $n\ge8$, are algebraic in local operators and have kernel decay uniform in $n$: $G_\sigma$, $Q_\sigma G_\sigma Q_\sigma^*$, $\Gamma_\sigma$, $H_\sigma$, $C_\sigma$, $\cD_\sigma$, $M$, $M^{-1}$, $P_\cV$ and $R$ (Sections~\ref{sec:scalar}, \ref{sec:scalar-avg}), $\bL_a$, $\bG$, $\bQ\bG\bQ^*$, $\boldsymbol\Gamma$, $\bH$, $\bC$ and $\bD$ (Sections~\ref{sec:gauge-coercivity}--\ref{sec:gauge-min}), and their products with the local operators of Section~\ref{sec:setup}. We denote their limits in the sense of Lemma~\ref{lem:ol-folding}(c) by the same letters. The identities of Part~I between them hold for the limits, in $\mathscr P$ and on the fields of polynomial growth. In particular
\[
\begin{gathered}
\bL_a=d^*d+dRd^*+a\bQ^*\bQ,\qquad\bL_a\bG=\bG\bL_a=I,\qquad\bH=\bG\bQ^*\boldsymbol\Gamma,\\
\bQ\bH=I,\qquad Rd^*\bH=0,\qquad d^*d\,\bH=\bQ^*(\boldsymbol\Gamma-a),
\end{gathered}
\]
$\bQ d=\dl Q_0$, $Q_0Q_0^*=I$, $R=I-P_\cV$, $R^2=R$, $Q_0G_0R=0$, $-\Lap G_0R=R$, and $P_\cV1=1$ for the constant field $1$.
\end{corollary}

\begin{proof}\hasdetail{cor:ol-part-one}
The families are algebraic in local operators by their definitions: $G_\sigma=(-\Lap+aQ_\sigma^*Q_\sigma)^{-1}$, invertible on every torus by Corollary~\ref{cor:scalar-coercive}; $\Gamma_\sigma=(Q_\sigma G_\sigma Q_\sigma^*)^{-1}$ (Proposition~\ref{prop:QGQ}); $H_\sigma=G_\sigma Q_\sigma^*\Gamma_\sigma$; $C_\sigma=G_\sigma-H_\sigma Q_\sigma G_\sigma$; $\cD_\sigma=\Gamma_\sigma-a$; $M=Q_0G_0^2Q_0^*$, invertible by Proposition~\ref{prop:R}(a); $P_\cV=G_0Q_0^*M^{-1}Q_0G_0$; $R=I-P_\cV$; $\bL_a=d^*d+dRd^*+a\bQ^*\bQ$; $\bG=\bL_a^{-1}$ (Theorem~\ref{thm:gauge-coercive}); $\boldsymbol\Gamma=(\bQ\bG\bQ^*)^{-1}$ (Proposition~\ref{prop:QGQ-gauge}); $\bH=\bG\bQ^*\boldsymbol\Gamma$; $\bC=\bG-\bG dRd^*\bG-\bH\bQ\bG$; $\bD=\boldsymbol\Gamma-a$.

\emph{Decay.} Part~I proves decay in the sense of Definition~\ref{def:decay}, with constants independent of $n$, for $G_\sigma$ (Theorem~\ref{thm:scalar-green}(a)), $Q_\sigma G_\sigma Q_\sigma^*$ and $\Gamma_\sigma$ (Proposition~\ref{prop:QGQ}), $H_\sigma$ (Theorem~\ref{thm:H-scalar}(a)), $C_\sigma$ (Proposition~\ref{prop:C-scalar}(d)), $\cD_\sigma$ (Proposition~\ref{prop:eff-action-scalar}), $M^{-1}$ and $P_\cV$ (Proposition~\ref{prop:R}(e) and its proof), $\bG$ (Theorem~\ref{thm:gauge-green}(a)), $\bQ\bG\bQ^*$ and $\boldsymbol\Gamma$ (Proposition~\ref{prop:QGQ-gauge}), $\bH$ (Theorem~\ref{thm:H-gauge}(c)), $\bC$ (Theorem~\ref{thm:C-gauge}(c)) and $\bD$ (Proposition~\ref{prop:eff-action-gauge}). A local family has such decay, with constants depending on $\eta$ (Lemma~\ref{lem:decay-calculus}(d)), and products of operators with decay have decay (Lemma~\ref{lem:decay-calculus}(b)); this covers $M$, $R$, $\bL_a$ and the products with local operators. Decay in the sense of Definition~\ref{def:decay} gives kernel decay. Indeed, let $s$ lie in the support of $\one_y$ and $s'$ in that of $\one_{y'}$, so that $|y-y'|\ge|s-s'|_n-\sqrt d$. Testing $K$ with $\delta_{s'}e$ and $\delta_se'$, unit vectors $e,e'\in E$, gives $|K(s,s')|\le(w(s')/w(s))^{1/2}\|\one_yK\one_{y'}\|\le\eta^{-d/2}Ce^{\gamma\sqrt d}e^{-\gamma|s-s'|_n}$.

\emph{Identities.} On every torus they are proved in Part~I: Theorem~\ref{thm:H-gauge}(a) for $\bQ\bH=I$ and $Rd^*\bH=0$; then $d^*d\bH=\bL_a\bH-dRd^*\bH-a\bQ^*\bQ\bH=\bQ^*\boldsymbol\Gamma-a\bQ^*$; Lemma~\ref{lem:avg}(b),(d) for $Q_0Q_0^*=I$ and $\bQ d=\dl Q_0$; Proposition~\ref{prop:R}(b),(d) for the identities of $R$. Lemma~\ref{lem:ol-folding}(d) carries them to the limits. The constant field $1$ is $n$-periodic, and $P_\cV1=1$ on every torus, since the constants lie in $\cV$ (Proposition~\ref{prop:R}(c)); so $P_\cV1=1$ for the limit, by Lemma~\ref{lem:ol-folding}(a),(c).
\end{proof}

\paragraph{Bloch fibres.} For $\theta\in\R^4$ let $\mathcal H_\theta(\Sigma)$ be the space of the fields $\varphi$ on $\Sigma$ with $\varphi(s+a)=e^{\ii\theta\cdot a}\varphi(s)$ for all $a\in\Z^4$, with the inner product $\ip\varphi\psi_\theta=\sum_{s\in\Sigma_0}w\,\ip{\varphi(s)}{\psi(s)}_E$. It depends on $\theta$ modulo $2\pi\Z^4$. Restriction to the cell identifies it unitarily with $\ell^2(\Sigma_0;E)$, with the weights, for every $\theta$. On the fine lattice, $\mathcal H_\theta$ is spanned by the plane waves of the momenta $\theta+2\pi l$, $l\in\Z^4/N\Z^4$, as in the Bloch decomposition of Section~\ref{sec:fourier}. By Lemma~\ref{lem:ol-folding}(a), every $K\in\mathscr P$ maps $\mathcal H_\theta(\Sigma)$ into $\mathcal H_\theta(\Sigma')$; the restriction is the \emph{fibre} $\hat K(\theta)$ of $K$. For a field $f$ on $\Sigma$ of finite support put $f_\theta(s)=\sum_{a\in\Z^4}e^{-\ii\theta\cdot a}f(s+a)$, and for an $n$-periodic field $f$ the same with the sum over $a\in\Z^4/n\Z^4$. Then $f_\theta\in\mathcal H_\theta(\Sigma)$, and $(Kf)_\theta=\hat K(\theta)f_\theta$, by the substitution $s'\mapsto s'-a$ in the sum defining $Kf$. Let $\Theta_n=(2\pi/n)\Z^4/2\pi\Z^4$.

\begin{lemma}[Bloch fibres]\label{lem:ol-bloch}
Let $K\in\mathscr P$, from the fields on $\Sigma$ to those on $\Sigma'$.
\begin{enumerate}[label=(\alph*)]
\item In the identification with the cells, $\hat K(\theta)$ has the matrix $\big(\sum_{a\in\Z^4}K(s,s'+a)e^{\ii\theta\cdot a}\big)_{s\in\Sigma'_0,\,s'\in\Sigma_0}$. It is real analytic and $2\pi\Z^4$-periodic in $\theta$. The fibres of sums, products and adjoints are the sums, products and adjoints of the fibres, and $\hat I(\theta)=I$.
\item \emph{(Tori.)} For $n\ge1$, $f\mapsto(f_\theta)_{\theta\in\Theta_n}$ maps the $n$-periodic fields onto $\bigoplus_{\theta\in\Theta_n}\mathcal H_\theta(\Sigma)$, with $f=n^{-d}\sum_\theta f_\theta$ and $\ip fg=n^{-d}\sum_\theta\ip{f_\theta}{g_\theta}_\theta$, and $[K]_n=\bigoplus_{\theta\in\Theta_n}\hat K(\theta)$.
\item \emph{(Infinite volume.)} The map $f\mapsto(f_\theta)$ extends to a unitary map from $\ell^2(\Sigma)$ onto the square integrable functions on $[-\pi,\pi]^d$ with values in $\ell^2(\Sigma_0;E)$, for the measure $(2\pi)^{-d}d\theta$. It carries $K$ to the multiplication by $\hat K(\theta)$. Hence $\|K\|=\max_\theta\|\hat K(\theta)\|$ on $\ell^2$. If $K=K^*$ and $c\in\R$, then $\ip f{Kf}\ge c\|f\|^2$ for all $f\in\ell^2(\Sigma)$ if and only if $\hat K(\theta)\ge c$ for all $\theta$.
\item \emph{(Inverses.)} If $\hat K(\theta)$ is invertible for every $\theta$, then $K$ has an inverse in $\mathscr P$, with the fibres $\hat K(\theta)^{-1}$. This holds if $K=K^*\ge c>0$ on $\ell^2(\Sigma)$.
\item \emph{(Transfer.)} Let $\mathfrak N$ be an infinite set of sides. Then $\|K\|=\sup_{n\in\mathfrak N}\|[K]_n\|$, and, if $K=K^*$, $K\ge c$ on $\ell^2(\Sigma)$ if and only if $[K]_n\ge c$ for all $n\in\mathfrak N$.
\item \emph{(Powers of $\langle\partial\rangle$.)} Let $\langle\partial\rangle=(1-\Lap)^{1/2}$ on the fine scalar or gauge fields, componentwise. For $t\in\R$, $\langle\partial\rangle^t\in\mathscr P$, its fibres are the powers $t/2$ of those of $1-\Lap$, and $[\langle\partial\rangle^t]_n=(1-\Lap)^{t/2}$ on the torus of side $n$.
\item \emph{(Part~I in infinite volume.)} The bounds of Part~I of the forms $X\ge c$ and $\|X\|\le C$, for operators $X$ built from those of Corollary~\ref{cor:ol-part-one} and from powers of $\langle\partial\rangle$ by sums and products, hold for the limits on $\ell^2$, with the same constants. For instance, with the constants of Corollary~\ref{cor:scalar-coercive}, Theorem~\ref{thm:gauge-coercive}, Propositions~\ref{prop:QGQ}, \ref{prop:QGQ-gauge} and~\ref{prop:R}(a),
\[
\begin{gathered}
-\Lap+aQ_\sigma^*Q_\sigma\ge c_0\langle\partial\rangle^2,\qquad\bL_a\ge c_g\langle\partial\rangle^2,\qquad\|\langle\partial\rangle\bG\langle\partial\rangle\|\le c_g^{-1},\\
c_{\rm QGQ}\le\bQ\bG\bQ^*\le2c_g^{-1},\qquad M\ge c_{\rm QGQ}^2 .
\end{gathered}
\]
Here $c_{\rm QGQ}$ is the constant $c_F$ of Corollary~\ref{cor:QKQ-lower}, renamed in this part to avoid a clash with the constant $c_{\rm F}$ of ($\Gamma$4).
\end{enumerate}
\end{lemma}

\begin{proof}\hasdetail{lem:ol-bloch}
(a) Every $s'\in\Sigma$ is $s''+a$ with unique $s''\in\Sigma_0$ and $a\in\Z^4$, and $\varphi(s''+a)=e^{\ii\theta\cdot a}\varphi(s'')$ for $\varphi\in\mathcal H_\theta$; this gives the matrix. For $s\in\Sigma'_0$ and $s'\in\Sigma_0$, $|x(s)-x(s')-a|\ge|a|-\sqrt d$, so the series converges absolutely and uniformly for $|\operatorname{Im}\theta|\le\gamma'<\gamma$ if $K$ has kernel decay $(C,\gamma)$; the matrix is holomorphic there. Sums and products of restrictions are restrictions. For adjoints, the grouping of the proof of Lemma~\ref{lem:ol-folding}(a), now modulo the diagonal action of all of $\Z^4$, gives $\ip\varphi{K\psi}_\theta=\ip{K^*\varphi}\psi_\theta$ for $\varphi\in\mathcal H_\theta(\Sigma')$, $\psi\in\mathcal H_\theta(\Sigma)$; the phases cancel.

(b) For $n$-periodic $f$, $n^{-d}\sum_{\theta\in\Theta_n}f_\theta(s)=\sum_{a}f(s+a)\,n^{-d}\sum_\theta e^{-\ii\theta\cdot a}=f(s)$. For $\theta,\theta'\in\Theta_n$, the sum over the sites of the torus of $w\ip{f_\theta}{g_{\theta'}}_E$ is $n^d\delta_{\theta\theta'}\ip{f_\theta}{g_\theta}_\theta$, which gives the inner product, and the dimensions agree. The last statement is $(Kf)_\theta=\hat K(\theta)f_\theta$ together with Lemma~\ref{lem:ol-folding}(a).

(c) For $f$ of finite support, $\int\|f_\theta\|_\theta^2\,(2\pi)^{-d}d\theta=\sum_{s\in\Sigma_0}w\sum_a|f(s+a)|^2=\|f\|^2$, by Parseval's identity for the Fourier series $\theta\mapsto f_\theta(s)$. The map is onto, since every square integrable $\ell^2(\Sigma_0;E)$-valued function on $[-\pi,\pi]^d$ is the sum of its Fourier series. The operator $K$ is bounded on $\ell^2$ (Schur's test), and $(Kf)_\theta=\hat K(\theta)f_\theta$ for $f$ of finite support, hence for all $f\in\ell^2(\Sigma)$. A multiplication by a continuous function of $\theta$ has the norm $\max_\theta\|\hat K(\theta)\|$, and it is $\ge c$ if and only if $\hat K(\theta)\ge c$ for all $\theta$.

(d) By (a), $\det\hat K(\theta)$ is holomorphic for $|\operatorname{Im}\theta|<\gamma$ and $2\pi\Z^4$-periodic, and it does not vanish for real $\theta$. By compactness it does not vanish for $|\operatorname{Im}\theta_\mu|\le\gamma'$, $\mu=1,\dots,d$, for some $\gamma'>0$, and there $G(\theta)=\hat K(\theta)^{-1}$ is holomorphic, periodic and bounded by some $C'$. Put $K'(s',s+a)=\int_{[-\pi,\pi]^d}G(\theta)_{s's}e^{-\ii\theta\cdot a}(2\pi)^{-d}d\theta$ for $s'\in\Sigma_0$, $s\in\Sigma'_0$, $a\in\Z^4$, and extend it periodically. Shifting the integration in $\theta_\mu$ to $\theta_\mu-\ii\gamma'\operatorname{sign}(a_\mu)$, which is allowed by the periodicity, gives $|K'(s',s+a)|\le C'e^{-\gamma'|a|_1}$, so $K'\in\mathscr P$. Its fibres are $G(\theta)$, by (a) and the inversion of Fourier series. So $KK'$ and $K'K$ have the fibres $I$. A kernel in $\mathscr P$ is determined by its fibres, its entries being their Fourier coefficients; so $KK'=I$ and $K'K=I$. If $K=K^*\ge c>0$, the fibres are $\ge c$ by (c), hence invertible.

(e) For every $\theta$ and $n$ there is $\theta'\in\Theta_n$ with $|\theta-\theta'|\le\pi\sqrt d/n$, so $\bigcup_{n\in\mathfrak N}\Theta_n$ is dense. The claims follow from (b), (c) and the continuity of $\hat K$.

(f) In the Fourier variables $q\in[-\pi/\eta,\pi/\eta]^d$ of $\eta\Z^4$, $1-\Lap$ is the multiplication by $\varrho(q)=1+\sum_\mu4\eta^{-2}\sin^2(\eta q_\mu/2)$ (Section~\ref{sec:fourier}). For $q_\mu=u_\mu+\ii v_\mu$ with $|v_\mu|\le\frac14$,
\[
\operatorname{Re}4\eta^{-2}\sin^2(\eta q_\mu/2)=2\eta^{-2}\big(1-\cos(\eta u_\mu)\cosh(\eta v_\mu)\big)\ge-2\eta^{-2}\big(\cosh(\eta v_\mu)-1\big)\ge-\tfrac1{15},
\]
since $\eta\le\frac13$. So $\operatorname{Re}\varrho\ge\frac12$ on this strip, and $\varrho^{t/2}$, with the principal branch, is holomorphic and bounded there, and $2\pi\eta^{-1}$-periodic in each $q_\mu$. The kernel of $\langle\partial\rangle^t$ is $(\eta/2\pi)^d\int e^{\ii q\cdot(x-x')}\varrho(q)^{t/2}dq$. Shifting the integration in $q_\mu$ by $\frac\ii4\operatorname{sign}(x-x')_\mu$ gives the bound $C_{\eta,t}e^{-|x-x'|_1/4}$. The kernel is invariant under the translations of $\eta\Z^4$, so $\langle\partial\rangle^t\in\mathscr P$. By (c), the unitary map of (c) carries $\langle\partial\rangle^t=(1-\Lap)^{t/2}$ both to the multiplication by its fibres and to the multiplication by the powers $t/2$ of the fibres of $1-\Lap$. Both are continuous in $\theta$, so they agree, and (b) gives the statement on the tori.

(g) By Corollary~\ref{cor:ol-part-one}, Lemma~\ref{lem:ol-folding}(a) and part~(f), the periodizations of such an $X$ are the corresponding operators on the tori of the sides $n\ge8$, where Part~I proves the bounds with constants independent of $n$. Part~(e) carries them to $\ell^2$. In the examples, $\ip f{(-\Lap+aQ_\sigma^*Q_\sigma)f}\ge c_0(\|f\|^2+\|\partial f\|^2)=c_0\ip f{(1-\Lap)f}$, and similarly for $\bL_a$. On every torus $\bL_a\ge c_g\langle\partial\rangle^2$ gives $\langle\partial\rangle^{-1}\bL_a\langle\partial\rangle^{-1}\ge c_g$, whose inverse is $\langle\partial\rangle\bG\langle\partial\rangle\le c_g^{-1}$.
\end{proof}

\paragraph{Trace per cell.} Let $\mathscr B$ be the set of the bounded operators $O$ from $\ell^2(\Sigma)$ to $\ell^2(\Sigma')$ of the form $(Of)_\theta=\hat O(\theta)f_\theta$, with a continuous $2\pi\Z^4$-periodic map $\theta\mapsto\hat O(\theta)$ into the linear maps from $\ell^2(\Sigma_0;E)$ to $\ell^2(\Sigma'_0;E')$. By Lemma~\ref{lem:ol-bloch}(c), $\mathscr P\subset\mathscr B$. The set $\mathscr B$ is closed under sums, products and adjoints, and for $O\in\mathscr B$ and $r>0$, $|O|^r=(O^*O)^{r/2}$ lies in $\mathscr B$, with the fibres $|\hat O(\theta)|^r$. Put $[O]_n=\bigoplus_{\theta\in\Theta_n}\hat O(\theta)$, as in Lemma~\ref{lem:ol-bloch}(b). For $O$ from $\Sigma$ to $\Sigma$ the \emph{trace per cell} is
\[
\operatorname{Tr}_1O=\int_{[-\pi,\pi]^d}\operatorname{tr}\hat O(\theta)\,\frac{d\theta}{(2\pi)^d},
\]
and for $O\in\mathscr B$ we put $\|O\|_r=(\operatorname{Tr}_1|O|^r)^{1/r}$, $1\le r<\infty$, and $\|O\|_\infty=\max_\theta\|\hat O(\theta)\|$, the norm of $O$ on $\ell^2$. These are the objects of Section~\ref{sec:ol-vacpol}. For $O\in\mathscr P$ with kernel decay $(C,\gamma)$ and $p\in\C^d$ with $|\operatorname{Im}p|<\gamma$, let $O^{(p)}$ be the kernel $e^{-\ii p\cdot(x(s)-x(s'))}O(s,s')$, and let $[\hat x_\mu,O]$ be the kernel $(x(s)-x(s'))_\mu O(s,s')$, $\mu=1,\dots,d$; put $[O,\hat x_\mu]=-[\hat x_\mu,O]$. On the fields of finite support, $O^{(p)}=e^{-\ii p\cdot\hat x}Oe^{\ii p\cdot\hat x}$ and $[\hat x_\mu,O]=\hat x_\mu O-O\hat x_\mu$, with the position operator $(\hat xf)(s)=x(s)f(s)$.

\begin{lemma}[Trace per cell]\label{lem:ol-trace-tools}
\begin{enumerate}[label=(\alph*)]
\item For $O\in\mathscr P$ from $\Sigma$ to $\Sigma$, $\operatorname{Tr}_1O=\sum_{s\in\Sigma_0}\operatorname{tr}_EO(s,s)$, and, if $O$ has kernel decay $(C,\gamma)$,
\[
\big|n^{-d}\operatorname{tr}[O]_n-\operatorname{Tr}_1O\big|\le C|\Sigma_0|(\dim E)S_{\gamma/2}e^{-\gamma n/2}\qquad(n\ge1).
\]
For $O\in\mathscr B$, $n^{-d}\operatorname{tr}[O]_n\to\operatorname{Tr}_1O$ as $n\to\infty$. Moreover $\operatorname{Tr}_1(O_1O_2)=\operatorname{Tr}_1(O_2O_1)$ for $O_1,O_2\in\mathscr B$ for which both products are defined.
\item \emph{(H\"older.)} Let $O_1,\dots,O_m\in\mathscr B$, with $O_1\cdots O_m$ defined from $\Sigma$ to $\Sigma$, and let $r_1,\dots,r_m\in[1,\infty]$ with $\sum_i1/r_i=1$. Then $|\operatorname{Tr}_1(O_1\cdots O_m)|\le\prod_i\|O_i\|_{r_i}$.
\item \emph{(Rank.)} If $\operatorname{rank}\hat Z(\theta)\le m$ for all $\theta$, then $\|Z\|_1\le m\|Z\|_\infty$. This holds with $m=|\Sigma''_0|\dim E''$ if $Z=Z_1Z_2$ with $Z_2\in\mathscr B$ into the fields on $\Sigma''$ with values in $E''$. For the scalar, respectively gauge, fields on $\Z^4$, $m=\dim E''$, respectively $m=d\dim E''$.
\item \emph{(Commutators with $\hat x$.)} Let $O\in\mathscr P$ have kernel decay $(C,\gamma)$. For $|\operatorname{Im}p|<\gamma$, $O^{(p)}$ and $[\hat x_\mu,O]$ lie in $\mathscr P$. The map $p\mapsto O^{(p)}$ is holomorphic in the norm of operators on $\ell^2$, and
\[
\partial_{p_\mu}O^{(p)}=-\ii\,[\hat x_\mu,O^{(p)}].
\]
Moreover $(O_1O_2)^{(p)}=O_1^{(p)}O_2^{(p)}$ and $[\hat x_\mu,O_1O_2]=[\hat x_\mu,O_1]O_2+O_1[\hat x_\mu,O_2]$; if $O^{-1}\in\mathscr P$, then $(O^{-1})^{(p)}=(O^{(p)})^{-1}$ for $|\operatorname{Im}p|$ smaller than the decay rates of $O$ and $O^{-1}$, and $[\hat x_\mu,O^{-1}]=O^{-1}[O,\hat x_\mu]O^{-1}$. For real $p$, $(O^{(p)})^*=(O^*)^{(p)}$, and the fibre of $O^{(p)}$ at $\theta$ is unitarily equivalent to $\hat O(\theta+p)$; so $\|O^{(p)}\|_r=\|O\|_r$ for $1\le r\le\infty$, and $\operatorname{Tr}_1O^{(p)}=\operatorname{Tr}_1O$.
\end{enumerate}
\end{lemma}

\begin{proof}
(a) By Lemma~\ref{lem:ol-bloch}(a), $\operatorname{tr}\hat O(\theta)=\sum_{s\in\Sigma_0}\sum_a\operatorname{tr}_EO(s,s+a)e^{\ii\theta\cdot a}$, and the integral keeps the term $a=0$. Every site of the torus is the class of a unique $s+b$ with $s\in\Sigma_0$ and $b\in\Z^4/n\Z^4$, so, by periodicity, $\operatorname{tr}[O]_n=n^d\sum_{s\in\Sigma_0}\sum_{a\in n\Z^4}\operatorname{tr}_EO(s,s+a)$. For $s\in\Sigma_0$ and $0\ne a\in n\Z^4$, $|\operatorname{tr}_EO(s,s+a)|\le(\dim E)Ce^{-\gamma|a|}\le(\dim E)Ce^{-\gamma n/2}e^{-\gamma|a|/2}$, which gives the bound. For $O\in\mathscr B$, $n^{-d}\operatorname{tr}[O]_n=n^{-d}\sum_{\theta\in\Theta_n}\operatorname{tr}\hat O(\theta)$ is a Riemann sum of a continuous periodic function. The cyclicity holds fibrewise.

(b) The fibre of $O_1\cdots O_m$ is $\hat O_1(\theta)\cdots\hat O_m(\theta)$. By H\"older's inequality for the Schatten norms of linear maps between finite-dimensional Hilbert spaces, $|\operatorname{tr}(\hat O_1(\theta)\cdots\hat O_m(\theta))|\le\prod_ia_i(\theta)$, with $a_i(\theta)=(\operatorname{tr}|\hat O_i(\theta)|^{r_i})^{1/r_i}$ for $r_i<\infty$ and $a_i(\theta)=\|\hat O_i(\theta)\|$ for $r_i=\infty$. H\"older's inequality for the probability measure $(2\pi)^{-d}d\theta$ on $[-\pi,\pi]^d$ gives $\int\prod_ia_i\le\prod_i\|O_i\|_{r_i}$.

(c) The trace norm of a linear map of rank at most $m$ is at most $m$ times its norm, since it is the sum of at most $m$ singular values. Apply this fibrewise and integrate. If $Z=Z_1Z_2$, then $\hat Z(\theta)=\hat Z_1(\theta)\hat Z_2(\theta)$ factors through $\ell^2(\Sigma''_0;E'')$, of dimension $|\Sigma''_0|\dim E''$; on $\Z^4$ the cell contains one point, and $d$ positively oriented bonds.

(d) $O^{(p)}$ is periodic and has kernel decay $(C,\gamma-|\operatorname{Im}p|)$, and $[\hat x_\mu,O]$ is periodic and has kernel decay $(C_{\gamma'},\gamma')$ for every $\gamma'<\gamma$. Put $\Delta x=x(s)-x(s')$. The kernel of $O^{(p+h)}-O^{(p)}+\ii h\cdot[\hat x,O^{(p)}]$ is $O(s,s')e^{-\ii p\cdot\Delta x}\big(e^{-\ii h\cdot\Delta x}-1+\ii h\cdot\Delta x\big)$, bounded by $C|h|^2|\Delta x|^2e^{-(\gamma-|\operatorname{Im}p|-|h|)|\Delta x|}$. By Schur's test its norm on $\ell^2$ is $O(|h|^2)$. This gives the derivative; the derivative is again of the same form, so the map is holomorphic. The product rules follow from $\Delta x=(x(s)-x(s'))+(x(s')-x(s''))$ in the kernel of the product. If $O^{-1}\in\mathscr P$, then $O^{(p)}(O^{-1})^{(p)}=I^{(p)}=I$, and the product rule applied to $OO^{-1}=I$ gives $[\hat x_\mu,O^{-1}]=-O^{-1}[\hat x_\mu,O]O^{-1}$. For real $p$ the identity for the adjoint is read off the kernels. For $\varphi\in\mathcal H_\theta$, the field $U_p\varphi=e^{\ii p\cdot x}\varphi$ lies in $\mathcal H_{\theta+p}$, and $U_p$ is unitary from $\mathcal H_\theta$ onto $\mathcal H_{\theta+p}$. Since $O^{(p)}\varphi=U_p^{-1}OU_p\varphi$, the fibre of $O^{(p)}$ at $\theta$ is $U_p^{-1}\hat O(\theta+p)U_p$, and that of $|O^{(p)}|^r$ is $U_p^{-1}|\hat O(\theta+p)|^rU_p$. The integrals over $\theta$ do not change under the shift by $p$, by periodicity.
\end{proof}

\begin{lemma}[Uniqueness among fields of polynomial growth]\label{lem:ol-liouville}
Let $K$ and $K'$ be kernels with kernel decay, from the fields on $\Sigma$ to those on $\Sigma'$ and back, with $K'K=I$. If $h$ is a field of polynomial growth on $\Sigma$ and $Kh=f$, then $h=K'f$; in particular, $Kh=0$ implies $h=0$. This applies to every $K\in\mathscr P$ with an inverse in $\mathscr P$, for instance to $\bL_a$, $\bQ\bG\bQ^*$, $-\Lap+aQ_\sigma^*Q_\sigma$ and $M$ (Corollary~\ref{cor:ol-part-one}).
\end{lemma}

\begin{proof}
By Lemma~\ref{lem:ol-folding}(a), $f$ has polynomial growth and $K'f=K'(Kh)=(K'K)h=h$.
\end{proof}

\begin{proposition}[The infinite-volume minimizer]\label{prop:ol-IM}
Let $\bH_\infty$ be the limit of the minimizer maps $\bH$ of Theorem~\ref{thm:E}, and let $\bG$, $\boldsymbol\Gamma$, $R$, $G_0$ be the limits of Corollary~\ref{cor:ol-part-one}; $\bQ$ and $Q_0$ are given on $\eta\Z^4$ by the formulas of Definition~\ref{def:averages}. The fields take values in $\mathfrak g^c$, and the operators act componentwise.
\begin{enumerate}[label=(\alph*)]
\item \emph{(Convergence.)} The kernels of $\bH$ on the tori of the sides $n\ge8$ are the periodizations of the kernel of $\bH_\infty$, an operator from the gauge fields on $\Z^4$ to those on $\eta\Z^4$. They have kernel decay $(C,\gamma)$ uniform in $n$, and they converge to it entrywise, with $|\bH(s,s')-\bH_\infty(s,s')|\le CS_{\gamma/4}e^{-\gamma n/2}$ for $|s-s'|\le n/4$.
\item \emph{(Uniqueness.)} Let $A$ be a gauge field on $\eta\Z^4$ of polynomial growth with $Rd^*A=0$ and $d^*dA=\bQ^*\omega$ for a gauge field $\omega$ on $\Z^4$ of polynomial growth. Then $A=\bH_\infty\bQ A$.
\item \emph{(Linear fields.)} For every curvature datum $F$ in $\mathfrak g^c$ and every $x_0\in\R^4$, $\bH_\infty B_F=A_F$, where $A_F$ and $B_F$ are the linear fields of $F$ with origin $x_0$ on the bonds of $\eta\Z^4$ and of $\Z^4$ (Section~\ref{subsec:gen-period}).
\item \emph{(Gradients.)} The family $\Psi=Q_0^*-G_0Rd^*dQ_0^*$, from the scalar fields on the unit torus to those on the fine torus, is algebraic in local operators, with kernel decay uniform in $n$. On every torus, and for the limits,
\[
Q_0\Psi=I,\qquad Rd^*d\,\Psi=0,\qquad\bH\dl=d\,\Psi .
\]
So $\bH_\infty$ maps the coarse gradient $\dl\lambda$ of every scalar field $\lambda$ on $\Z^4$ of polynomial growth to the fine gradient $d\Psi_\infty\lambda$. If $\varphi$ is a polynomial on $\R^4$ of degree at most $2$, restricted to $\eta\Z^4$, then $\bH_\infty\dl Q_0\varphi=d\varphi$.
\item \emph{(Coarse linear fields.)} Let $p,a\in\C^4$ and $T\in\mathfrak g^c$. For an oriented bond $b$ let $m_b$ be its midpoint and $t_b\in\{\pm e_1,\dots,\pm e_4\}$ its direction. Then $\bH_\infty$ maps the field $c\mapsto(p\cdot m_c)(a\cdot t_c)T$ on the bonds of $\Z^4$ to the field $b\mapsto(p\cdot m_b)(a\cdot t_b)T$ on the bonds of $\eta\Z^4$, and the constant field $c\mapsto(a\cdot t_c)T$ to the constant field $b\mapsto(a\cdot t_b)T$.
\end{enumerate}
\end{proposition}

\begin{proof}\hasdetail{prop:ol-IM}
(a) By Corollary~\ref{cor:ol-part-one} and Lemma~\ref{lem:ol-folding}(c).

(b) Put $B=\bQ A$ and $h=A-\bH_\infty B$. By Lemma~\ref{lem:ol-folding}(a), $B$ and $h$ have polynomial growth. By Corollary~\ref{cor:ol-part-one}, $\bQ\bH_\infty=I$, $Rd^*\bH_\infty=0$ and $d^*d\bH_\infty=\bQ^*(\boldsymbol\Gamma-a)$. So $\bQ h=0$, $Rd^*h=0$, and
\[
\bL_ah=d^*dh+dRd^*h+a\bQ^*\bQ h=\bQ^*\omega',\qquad\omega'=\omega-(\boldsymbol\Gamma-a)B,
\]
a field of polynomial growth. By Lemma~\ref{lem:ol-liouville} for $\bL_a$, with the inverse $\bG$, $h=\bG\bQ^*\omega'$. Then $0=\bQ h=(\bQ\bG\bQ^*)\omega'$, and Lemma~\ref{lem:ol-liouville} for $\bQ\bG\bQ^*$, with the inverse $\boldsymbol\Gamma$, gives $\omega'=0$. Hence $h=0$.

(c) $A_F$ has linear growth, and $\bQ A_F=B_F$ by Lemma~\ref{lem:ol-linear-average}(a). On the bond $b=(x,x+\eta e_\nu)$, $A_F(b)=\frac12\sum_\kappa(x+\frac12\eta e_\nu-x_0)_\kappa F_{\kappa\nu}$ does not depend on $x_\nu$, since $F_{\nu\nu}=0$. So $\partial_\nu^*A_{F,\nu}=0$ for every $\nu$, $d^*A_F=0$, and $Rd^*A_F=0$. Moreover $(dA_F)_{\mu\nu}=\partial_\mu A_{F,\nu}-\partial_\nu A_{F,\mu}=F_{\mu\nu}$ is constant (Section~\ref{subsec:gen-period}), so $d^*dA_F=0$. Part~(b) with $\omega=0$ gives $A_F=\bH_\infty B_F$.

(d) Let $n\ge8$. By Lemma~\ref{lem:avg}(b), $Q_0Q_0^*=I$, and by Proposition~\ref{prop:R}(b),(d), $Q_0G_0h=0$ and $-\Lap G_0h=h$ for $h$ in the range of $R$. With $d^*d=-\Lap$ on scalar fields,
\[
\begin{aligned}
Q_0\Psi&=I-Q_0G_0R(-\Lap Q_0^*)=I,\\
Rd^*d\Psi&=R(-\Lap)Q_0^*-R(-\Lap G_0R)(-\Lap)Q_0^*=R(-\Lap)Q_0^*-R^2(-\Lap)Q_0^*=0 .
\end{aligned}
\]
Let $\lambda$ be a scalar field on the unit torus and $A=d\Psi\lambda$. Then $\bQ A=\dl Q_0\Psi\lambda=\dl\lambda$ by Lemma~\ref{lem:avg}(d), $Rd^*A=0$ and $dA=0$. By Theorem~\ref{thm:H-gauge}(b), $\bH\dl\lambda$ is the unique minimizer of $\|dA\|^2$ on $\{A:\bQ A=\dl\lambda,\ Rd^*A=0\}$, and $d\Psi\lambda$ attains the value $0$; so $\bH\dl\lambda=d\Psi\lambda$. The family $\Psi$ is algebraic in local operators, with kernel decay uniform in $n$, by Corollary~\ref{cor:ol-part-one} and Lemma~\ref{lem:ol-folding}(d), which also carries the three identities to the limits, and Lemma~\ref{lem:ol-folding}(a) applies them to the fields of polynomial growth.

Let $\varphi$ be a polynomial of degree at most $2$, and $A=d\varphi$. Then $A$ has polynomial growth and $dA=0$. On the lattice, $-\Lap x_\mu^2=-2$, and $-\Lap$ annihilates $x_\mu x_\nu$ for $\mu\ne\nu$, $x_\mu$ and $1$; so $d^*A=-\Lap\varphi$ is constant. By Corollary~\ref{cor:ol-part-one}, $R1=1-P_\cV1=0$, so $Rd^*A=0$. Part~(b) with $\omega=0$ and $\bQ d\varphi=\dl Q_0\varphi$ give $d\varphi=\bH_\infty\dl Q_0\varphi$.

(e) Put $\varphi(x)=\frac12(p\cdot x)(a\cdot x)$ and let $F$ be the curvature datum $F_{\kappa\lambda}=(p_\kappa a_\lambda-p_\lambda a_\kappa)T$, with origin $0$. For an oriented bond $b$ of $\eta\Z^4$, $b_\pm=m_b\pm\frac12\eta t_b$, so
\[
(d\varphi)(b)=\eta^{-1}\big(\varphi(b_+)-\varphi(b_-)\big)=\tfrac12\big((p\cdot m_b)(a\cdot t_b)+(a\cdot m_b)(p\cdot t_b)\big),
\]
exactly, while $A_F(b)=\frac12m_b^{\sf T}Ft_b=\frac12\big((p\cdot m_b)(a\cdot t_b)-(a\cdot m_b)(p\cdot t_b)\big)T$. Hence $(p\cdot m_b)(a\cdot t_b)T=A_F(b)+(d\varphi)(b)T$. The same computation on $\Z^4$ gives $(p\cdot m_c)(a\cdot t_c)T=B_F(c)+(\dl\varphi)(c)T$, with $\varphi$ restricted to $\Z^4$. The blocks are centred, so the average of $\varphi(y+t)$ over the points $t$ of $B(0)$ is $\varphi(y)$ plus the constant $\frac12\sum_\mu p_\mu a_\mu\,\eta^d\sum_{t\in B(0)}t_\mu^2$; hence $\dl Q_0\varphi=\dl\varphi$ on $\Z^4$. By (c) and (d), $\bH_\infty$ maps $B_F+\dl\varphi\,T$ to $A_F+d\varphi\,T$. For the constant fields, take $\varphi(x)=a\cdot x$: then $Q_0\varphi=\varphi$ on $\Z^4$, $(\dl\varphi)(c)=a\cdot t_c$ and $(d\varphi)(b)=a\cdot t_b$, and (d) applies.
\end{proof}

In particular $\bH_\infty$ maps linear fields to linear fields exactly, not only up to gauge. For the Bloch components this gives the following.

\begin{corollary}[Umklapp components]\label{cor:ol-umklapp}
Let $\delta^{-1}$ be a divisor of $N$, and let $\mathbf K\in\mathscr P$ map the gauge fields on $\Z^4$ to the gauge fields on $\delta\Z^4$, with real kernel, such that for all $p,a\in\C^4$ it maps the field $c\mapsto a\cdot t_c$ to $b\mapsto a\cdot t_b$, and $c\mapsto(p\cdot m_c)(a\cdot t_c)$ to $b\mapsto(p\cdot m_b)(a\cdot t_b)$. For $q\in\R^4$ and $\varepsilon\in\C^4$ let $e_{q,\varepsilon}(c)=e^{\ii q\cdot m_c}(\varepsilon\cdot t_c)$ on the bonds of $\Z^4$, and $e'_{q,\varepsilon}$ the same field on the bonds of $\delta\Z^4$. Let $\mathcal L=\{0,\dots,\delta^{-1}-1\}^4$. Then, for $p\in\R^4$,
\[
\mathbf Ke_{p,\varepsilon}=\sum_{\ell\in\mathcal L}e'_{p+2\pi\ell,\,k_\ell(p)\varepsilon},
\]
a finite sum, with $d\times d$ matrices $k_\ell(p)$ that are real analytic in $p$, and $k_0(p)=I+O(|p|^2)$, $k_\ell(p)=O(|p|^2)$ for $\ell\ne0$. By Proposition~\ref{prop:ol-IM}(e), this applies to $\bH_\infty$, with $\delta=\eta$.
\end{corollary}

\begin{proof}\hasdetail{cor:ol-umklapp}
Fix a direction $\mu$. The functions $m\mapsto e^{\ii(p+2\pi\ell)\cdot m}$, $\ell\in\mathcal L$, on the midpoints of the bonds of $\delta\Z^4$ of direction $\mu$ are $p$-quasi-periodic under $\Z^4$. On the midpoints $\delta(z+\frac12e_\mu)$ of the bonds of direction $\mu$ of a cell, with $z$ running over a set of representatives of $\Z^4/\delta^{-1}\Z^4$, they are $e^{\ii p\cdot\delta z}$ times the characters $z\mapsto e^{2\pi\ii\delta\ell\cdot z}$ of $\Z^4/\delta^{-1}\Z^4$, up to nonzero constants, so they are linearly independent and span the $p$-quasi-periodic functions. Since $\mathbf K$ is periodic, $\mathbf Ke_{p,\varepsilon}$ is $p$-quasi-periodic, which gives the expansion. The coefficients are linear combinations of the values of $f_p=e^{-\ii p\cdot m_b}(\mathbf Ke_{p,\varepsilon})(b)=\sum_c\mathbf K(b,c)e^{\ii p\cdot(m_c-m_b)}(\varepsilon\cdot t_c)$ on one cell, which are holomorphic in $p$ near $\R^4$, by the kernel decay. Write $\chi_\ell(b)=e^{2\pi\ii\ell\cdot m_b}$, so that $f_p=\sum_\ell\chi_\ell\,(k_\ell(p)\varepsilon)\cdot t_b$, the $\chi_\ell$ being linearly independent. At $p=0$, $f_0=\mathbf Ke_{0,\varepsilon}=\varepsilon\cdot t_b$, so $k_\ell(0)=\delta_{\ell0}I$. Differentiating under the sum, which the kernel decay allows,
\[
\frac d{d\tau}\Big|_{\tau=0}f_{\tau p}(b)=\ii\,\mathbf K\big[(p\cdot m_c)(\varepsilon\cdot t_c)\big](b)-\ii(p\cdot m_b)\,\mathbf K\big[\varepsilon\cdot t_c\big](b)=\ii(p\cdot m_b)(\varepsilon\cdot t_b)-\ii(p\cdot m_b)(\varepsilon\cdot t_b)=0 .
\]
So $\sum_\ell\chi_\ell\,\big((p\cdot\nabla)k_\ell(0)\varepsilon\big)\cdot t_b=0$, and $(p\cdot\nabla)k_\ell(0)=0$ for all $\ell$ and $p$.
\end{proof}

\subsection{The change of scale}\label{subsec:ol-scale}

For $j<k$ and a coarse configuration $V$ of scale $k$, let $V_{j+1}(V)=U_k(V)^{(j+1)}$, the configuration of scale $j+1$ determined by the background. In the units of step $j+1$ it lives on a unit lattice, and $V\mapsto V_{j+1}(V)$ is real analytic and covariant. Put $N'=L^{k-j-1}$. In the units of scale $k$, the configurations of scale $j+1$ live on $\Lam'=N'^{-1}\Z^4$, and a configuration $e^{\ii N'^{-1}B'}$ there has the Lie field $B'$, in the convention of Part~IV. Let $\bQ^{\rm pl}_{j+1}$ be the linearization at $1$ of the plain average $U\mapsto U^{(j+1)}$, from the fine lattice to $\Lam'$, in the units of scale $k$: $(e^{\ii\eta tA})^{(j+1)}=e^{\ii N'^{-1}t\bQ^{\rm pl}_{j+1}A+O(t^2)}$. It is Lemma~\ref{lem:ol-linear-average}(b) with $j+1$ in place of $k$, read in the units of scale $k$. In the units of scale $j+1$ the positions are multiplied by $N'$, and the Lie field of this configuration is $N'^{-1}t\bQ^{\rm pl}_{j+1}A+O(t^2)$. The kernel of a second variation, a function of pairs of bonds of $\Lam'$, does not change with the units.

\begin{lemma}[Change of scale]\label{lem:ol-pullback}
Let $\Phi$ be as above, a function of configurations of scale $j+1$, with $\Pi_\Phi(p)=O(|p|^2)$, and put $\Psi(V)=\Phi(V_{j+1}(V))$. Then $\Psi$ satisfies the assumptions made above Definition~\ref{def:marginal-tensor}, its infinite-volume second variation is
\[
\Pi_\Psi=N'^{-2}\,\mathbf K^{\sf T}\Pi_\Phi\mathbf K,\qquad\mathbf K=\bQ^{\rm pl}_{j+1}\bH_\infty ,
\]
$\Pi_\Psi(p)=O(|p|^2)$, and $T_\Psi=T_\Phi$, each in the units of its own scale.
\end{lemma}

\begin{proof}\hasdetail{lem:ol-pullback}
\emph{The linearization.} Let $V=e^{\ii B}$ with $B$ small, on a torus. The reference $U_0=1$ is an element of the minimal orbit of the configuration $1$, so Corollary~\ref{cor:analytic} applies, with $J_{U_0}=0$ and $W_k=1$: the minimal orbit of $V$ contains $\mathsf U(A_*(B))$, a gauge transform of $e^{\ii\eta A_*(B)}$, and $A_*$ is analytic in $B$. By Theorem~\ref{thm:variational}(c), $A_*(B)=\bH B+O(|B|^2)$, with the minimizer map $\bH$ of Part~I, since at $U_0=1$ the averaging pair is $(Q_0,\bQ)$ (Proposition~\ref{prop:linearized-averages}(c)). Indeed, at $U_0=1$ the defect $E$ vanishes and $(Q'_k,\bQ_k)=(Q_0,\bQ)$ (Proposition~\ref{prop:linearized-averages}(b),(c)), so $\bQ_\pi=\bQ$ (proof of Theorem~\ref{thm:cov-min}(a)); the form of Theorem~\ref{thm:variational} has $W=0$ (Section~\ref{sec:var-solution}), so $\fA_\pi=\bL_a$ (Section~\ref{sec:cov-gauge}, after Lemma~\ref{lem:slice}); $P_\cV=G_0Q_0^*M^{-1}Q_0G_0$; and the minimizer map $\bG\bQ_\pi^T\boldsymbol\Gamma_\pi$ of Theorem~\ref{thm:cov-min}(b) is the map $\bH=\bG\bQ^*\boldsymbol\Gamma$ of Theorem~\ref{thm:E}. The configuration $U_k(V)$ is a gauge transform of $\mathsf U(A_*(B))$, the averaging is covariant under all real gauge transformations (Lemma~\ref{lem:one-step-basic}(b)), and $\Phi$ is gauge invariant. Hence $\Psi(V)=\Phi\big((e^{\ii\eta A_*(B)})^{(j+1)}\big)$. Since the first variation of $\Phi$ at $1$ vanishes, only the linearization of $B\mapsto(e^{\ii\eta A_*(B)})^{(j+1)}$ enters the second variation, and on every torus $\Pi_\Psi=N'^{-2}\mathbf K^{\sf T}\Pi_\Phi\mathbf K$ with $\mathbf K=\bQ^{\rm pl}_{j+1}\bH$; the factor $N'^{-2}$ is the change of units of the Lie fields.

\emph{Infinite volume.} The family $\bQ^{\rm pl}_{j+1}$ is local: $(\bQ^{\rm pl}_{j+1}A)(c')$ depends only on $A$ on the bonds of the blocks of scale $j+1$ of the end points of $c'$ (Lemma~\ref{lem:one-step-basic}(c)), it is given by one formula on every torus, and it commutes with the translations of $\Lam'$ (Lemma~\ref{lem:one-step-basic}(d)). By Proposition~\ref{prop:ol-IM}(a) and Lemma~\ref{lem:ol-folding}(a),(c), the maps $\mathbf K$ on the tori are the periodizations of $\mathbf K_\infty=\bQ^{\rm pl}_{j+1}\bH_\infty\in\mathscr P$, and they converge to it with uniform decay. By the assumption on $\Phi$ and Lemma~\ref{lem:ol-folding}(e), the second variations of $\Psi$ on the tori converge with uniform decay to $N'^{-2}\mathbf K_\infty^{\sf T}\Pi_\Phi\mathbf K_\infty$. This kernel is invariant under the translations of $\Z^4$, since $\mathbf K_\infty$ commutes with them and $\Pi_\Phi$ with those of $\Lam'\supset\Z^4$. The function $\Psi$ is real analytic and gauge invariant, as $\Phi$ is and $V_{j+1}$ is real analytic and covariant, and its first variation at $1$ vanishes. So $\Psi$ satisfies the assumptions above Definition~\ref{def:marginal-tensor}. From now on $\mathbf K=\mathbf K_\infty$.

\emph{The Umklapp components.} As in Definition~\ref{def:marginal-tensor}, we suppress the Lie algebra indices and use the bond-midpoint convention, with the fields $e_{q,\varepsilon}$ on $\Z^4$ and $e'_{q,\varepsilon}$ on $\Lam'$ of Corollary~\ref{cor:ol-umklapp}. The map $\mathbf K$ maps the constant fields and the fields $c\mapsto(p\cdot m_c)(a\cdot t_c)$ on $\Z^4$ to the same fields on $\Lam'$. For $\bH_\infty$, with $\eta\Z^4$ in place of $\Lam'$, this is Proposition~\ref{prop:ol-IM}(e). For $\bQ^{\rm pl}_{j+1}$, write the fine field as $A_F+d\varphi$, with $F=p\wedge a$ and $\varphi(x)=\frac12(p\cdot x)(a\cdot x)$, as in the proof of Proposition~\ref{prop:ol-IM}(e). By Lemma~\ref{lem:ol-linear-average}(b), $\bQ^{\rm pl}_{j+1}A_F$ is the linear field of $F$ on $\Lam'$, with no gradient, and by the gauge covariance in its proof $\bQ^{\rm pl}_{j+1}d\varphi=d'\varphi$, the lattice gradient of $\varphi$ on $\Lam'$, $(d'\varphi)(c')=N'\big(\varphi(c'_+)-\varphi(c'_-)\big)$. Their sum is the field $(p\cdot m)(a\cdot t)$ on $\Lam'$, by the same identity on $\Lam'$; constant fields are the case $\varphi(x)=a\cdot x$. So Corollary~\ref{cor:ol-umklapp} applies, with $\delta=N'^{-1}$: $\mathbf Ke_{p,\varepsilon}=\sum_{\ell\in\mathcal L}e'_{p+2\pi\ell,k_\ell(p)\varepsilon}$, a finite sum over $N'^d$ Umklapp momenta, with $k_0(p)=I+O(|p|^2)$ and $k_\ell(p)=O(|p|^2)$ for $\ell\ne0$. Since $\mathbf K$ is real, $\mathbf Ke_{-p,\varepsilon}=\sum_\ell e'_{-p-2\pi\ell,\overline{k_\ell(p)}\varepsilon}$.

\emph{The tensor.} For a $(-p)$-quasi-periodic field $g$ and a $p$-quasi-periodic field $f$ on the same bonds, put $(g,f)_1=\sum g(b)f(b)$, the sum over the positively oriented bonds with base point in a cell. For a periodic kernel $X$, $(g,Xf)_1=(X^{\sf T}g,f)_1$, by the grouping of the proof of Lemma~\ref{lem:ol-folding}(a). By the definition of the Fourier transform and by this identity, $\varepsilon^{\sf T}\Pi_\Psi(p)\varepsilon'=(e_{-p,\varepsilon},\Pi_\Psi e_{p,\varepsilon'})_1=N'^{-2}(\mathbf Ke_{-p,\varepsilon},\Pi_\Phi\mathbf Ke_{p,\varepsilon'})_1$, the last pairing over the bonds of $\Lam'$ with base point in a cell of $\Z^4$. The kernel $\Pi_\Phi$ commutes with the translations of $\Lam'$, so $\Pi_\Phi e'_{q,\varepsilon}=e'_{q,\Pi_\Phi(q/N')\varepsilon}$, with the Fourier transform $\Pi_\Phi(\cdot)$ in the units of scale $j+1$. On a cell of $\Z^4$, $\Lam'$ has $N'^d$ bonds of each direction, and $(e'_{-p-2\pi\ell,\varepsilon},e'_{p+2\pi\ell',\varepsilon'})_1=N'^d\delta_{\ell\ell'}\varepsilon^{\sf T}\varepsilon'$ for $\ell,\ell'\in\mathcal L$, by the orthogonality of the characters. Hence
\[
\Pi_\Psi(p)=N'^{d-2}\sum_{\ell\in\mathcal L}k_\ell(p)^*\,\Pi_\Phi\big((p+2\pi\ell)/N'\big)\,k_\ell(p).
\]
The term $\ell=0$ is $N'^{d-2}\Pi_\Phi(p/N')+O(|p|^4)=N'^{d-2}T_\Phi(p/N')+O(|p|^3)=N'^{d-4}T_\Phi(p)+O(|p|^3)$, since $\Pi_\Phi(q)=O(|q|^2)$ and $T_\Phi$ is homogeneous of degree $2$. For $\ell\ne0$ the momenta $(p+2\pi\ell)/N'$, $|p|\le1$, lie in a compact set, where $\Pi_\Phi$ is bounded, since its kernel decays exponentially; so these terms are $O(|p|^4)$. Hence $\Pi_\Psi(p)=N'^{d-4}T_\Phi(p)+O(|p|^3)$. In $d=4$ the factor is $1$.
\end{proof}

No gauge argument is needed for the Umklapp components: by Proposition~\ref{prop:ol-IM}(e) they vanish to second order in $p$. The proof uses the infinite-volume minimizer of Proposition~\ref{prop:ol-IM} and the averages of linear fields of Lemma~\ref{lem:ol-linear-average}, both exact.

\begin{corollary}[The cumulative one-loop tensor]\label{cor:ol-cumulative}
Let $T^{(1)}_{j+1}$ be the marginal tensor of $V'\mapsto\Gamma^{\rm step}_j(U_{j+1}(V'))$, a function of the configuration $V'$ of scale $j+1$. Assume that the second variations at $1$ of these functions and of $\Gamma^{(1)}_k$ on the tori converge with uniform exponential decay as the side tends to infinity (Lemma~\ref{lem:ol-folding}(e)). Then they satisfy the hypothesis of Definition~\ref{def:marginal-tensor}, each of their marginal tensors is a multiple of $\delta_{\mu\nu}|p|^2-p_\mu p_\nu$, and
\[
\sum_{j<k}T^{(1)}_{j+1}=T_{\Gamma^{(1)}_k}.
\]
\end{corollary}

The hypothesis on the second variations is proved in Lemma~\ref{lem:ol-torus-limit}(c) for $\Gamma^{(1)}_k$, and in Lemma~\ref{lem:ol-limits} for the step functions.

\begin{proof}\hasdetail{cor:ol-cumulative}
The hypothesis of Definition~\ref{def:marginal-tensor} and the form of the tensors are Lemma~\ref{lem:ol-symmetry}(c). By Proposition~\ref{prop:ol-composition} and Lemma~\ref{lem:composition}, $\Gamma^{(1)}_k(V)=\sum_{j<k}\Gamma^{\rm step}_j(U_{j+1}(V_{j+1}(V)))$ up to a constant. Apply Lemma~\ref{lem:ol-pullback} to each term; marginal tensors are additive.
\end{proof}

\paragraph{Relation to Ba\l aban's $\beta$-functions.} The averaging of Part~IV and the gauge fixing of Part~VII are Ba\l aban's \cite[(0.3)--(0.19)]{BalYM4I}. Ba\l aban reads off $\beta_{j+1}$ from the components with $\mu\ne\nu$ of the marginal tensor of the new generation, in infinite volume \cite[(1.20)--(1.22), (2.14)]{BalYM4I}, and its one-loop part from the term of order $g^0$ of the new generation. That term is $\Gamma^{\rm step}_j$, and by Lemma~\ref{lem:ol-symmetry} its marginal tensor is $T^{(1)}_{j+1}=\beta^{(1)}_{j+1}(\delta_{\mu\nu}|p|^2-p_\mu p_\nu)$. The flow of Part~X reads $\beta_{j+1}$ off the symmetrized new generation on the given torus, through its period density (Proposition~\ref{prop:period-iso}). The marginal parts of the old generations, which must be cancelled before $\beta$ is read off, are cancelled there by their counterterms, inside the old activities of the step (Definition~\ref{def:old-activities}); they do not enter the one-loop part. Section~\ref{sec:ol-uv} compares the two readings. The one-loop part of $\beta_{j+1}$ is $\beta^{(1)}_{j+1}$ up to $O(e^{-cn_{j+1}/M})$ ((F2) of Proposition~\ref{prop:flow-F}), the rest of $\beta_{j+1}$ is $O(g_j^2p_0(g_j)^6)$ ((F3)), and, at fixed couplings, $\beta_{j+1}$ converges as the torus grows to the coefficient that Ba\l aban's definition assigns to the generations of this note (Proposition~\ref{prop:beta-identification}, Section~\ref{sec:gen-locality}); no rate is proved for this convergence. The one-loop parts $\beta^{(1)}_{j+1}$, summed over the steps, are the coefficient of the marginal tensor of the single function $\Gamma^{(1)}_k$ (Corollary~\ref{cor:ol-cumulative} and Theorem~\ref{thm:ol-cumulative}), which the next section writes in the Landau chart.
